\documentclass[../main.tex]{subfiles}
\begin{document}
%==========================================TIÊU ĐỀ TẬP========================================
\begin{table}[h]
    \begin{adjustwidth}{-1cm}{}
        \begin{tabular}{>{\centering\arraybackslash}p{6cm}|>{\raggedright\arraybackslash}p{14cm}}
            \multirow{5}{6cm}
            % Cột bên trái: ÉPISODE và số tập
            {\\[-40pt]
                \flaregothic\fontsize{20pt}{20pt}\selectfont \centering
                \color{eptype}HISTOIRE\color{black}\\
                \brian\fontsize{72pt}{72pt}\selectfont
                \color{epnum}2
                \\[18pt]
                % \flaregothic\fontsize{14pt}{14pt}\selectfont
                % \color{epattr}BIENVENUE, \uline{\color{Banpire} Mel} !
                % \flaregothic\fontsize{18pt}{18pt}\selectfont
                % \color{epattr}ÉPISODE PERDU
                \color{black}
            }
             &
            % Dòng thứ nhất: tên tập tiếng Nhật
            {\fotskip\fontsize{18pt}{18pt}\selectfont \ruby{ショッピング}{しょっぴんぐ}\ruby{モール}{もーる}\ruby{大}{だい}\ruby{騒}{そう}\ruby{動}{どう}！} \\[6pt]
            % Dòng thứ hai: tên tập tiếng Anh
             & {\cafeteria\fontsize{18pt}{18pt}\selectfont Shopping Mall Mayhem!}                                                \\[4pt]
            % Dòng thứ ba: tên tập tiếng Việt
             & {\vnmsans\fontsize{18pt}{18pt}\selectfont Đại náo trung tâm mua sắm!}                                             \\[8pt]
            % Dòng thứ tư: Nhân vật xuất hiện
             & {\sen Track Used: Anniversary (アニバーサリー)}
        \end{tabular}
    \end{adjustwidth}
\end{table}
%===========================================NỘI DUNG==========================================
\color{black}

\begin{center}
    \uline{Trung tâm thương mại AEON - quận Ryuuhen.}
\end{center}

Nằm trong một trong những khu sầm uất nhất thành phố, AEON Ryuuhen là nơi nổi tiếng với sự đa dạng của các cửa hàng, từ thời trang, thực phẩm, đến các mặt hàng trang trí. Đây cũng chính là nơi mà trước đó, Miko, Mio, Korone, Lamy, Botan và Polka đã từng ghé qua để mua sắm vật dụng trang trí văn phòng hồi cuối đông\footnote{Xem lại {\flaregothic ÉPISODE 95}, quyển số Bảy.}. Nene cũng tới hồi đó, nhưng cô đến để theo dõi mọi người, đặc biệt là từ cậu Tổng, người mà cô gặp sau này ở công ty.

Lần này, họ quay lại cùng gần như toàn bộ thành viên \hll\ để chuẩn bị cho sinh nhật Kuon.

Mio nheo mắt nhìn xung quanh, rồi buông một tiếng thở dài: ``\dots Chúng ta lại quay trở lại chỗ này à. Kể từ lần mà chúng ta tới đây mua đồ để trang trí văn phòng cuối đông vừa rồi ha.''

Miko, với tinh thần hăng hái như mọi khi, giơ tay hô lớn: ``Nếu muốn mua đồ thì ở JUSCO cái gì cũng có, ne!''

Nàng Sói đen bật cười: ``Em biết là thế, nhưng lần trước chị đến đây với Korone rồi, cơ mà\dots''

Trước khi cô kịp nói hết câu, cô nhận ra một cảnh tượng quen thuộc: cả đám \hll\ đang chạy tán loạn trong trung tâm thương mại. Một số người lần đầu tiên đặt chân đến đây trông như những đứa trẻ lạc vào xứ sở thần tiên, mắt sáng rực vì choáng ngợp trước không gian rộng lớn. Một số khác thì hào hứng bàn luận về những món đồ trưng bày, sôi nổi chẳng khác gì một chuyến dã ngoại.

Nàng Sói đen chỉ có thể xoa trán thở dài: ``Mới đến thôi mà đã thế này rồi\dots''

Giữa không khí rộn ràng ấy, A-chan và Sora nhanh chóng đứng ra tập hợp mọi người lại. Nàng Thần tượng vẫy tay gọi: ``Tập trung lại đây nào mọi người ơi!'

Nhưng để gom ba mươi mốt con người đang hưng phấn thế này không phải chuyện dễ. Phải mất tận năm phút, cộng thêm vài tiếng hét từ Fubuki và một cú vỗ tay đầy uy quyền của Noel, mọi người mới chịu ổn định hàng ngũ.

Nàng ``Tăng ca'' đẩy gọng kính, lên tiếng: ``Bây giờ, theo ý tưởng của YAGOO-sama, chúng ta sẽ chia ra làm bốn nhóm, mỗi nhóm khoảng từ sáu đến bảy người.''

Sora tiếp lời: ``Một nhóm sẽ mua bánh và đồ ngọt, một nhóm sẽ mua đồ mặn, một nhóm chuẩn bị quà, và một nhóm mua đồ trang trí nha!''

Cả nhóm đồng thanh, chào theo kiểu nhà binh: ``\textbf{\eng{Rõ, thưa sếp!}}''

Và thế là, cả ba mươi mốt con người bắt đầu tản ra, mỗi nhóm mang một nhiệm vụ riêng, bước vào hành trình mua sắm cho sự kiện đặc biệt này.

\vspace{-40pt}
\subsubsection*{\phantomsection\label{P1.1a}}
\addcontentsline{toc}{subsubsection}{Đội \#{\digi \ 1F0972} - Aqua, Ayame, Shion, Korone và Okayu}

\begin{center}
    {\color{T1} Nhóm 1 - Aqua, Ayame, Shion, Korone và Okayu\\Nguyên liệu làm bánh}
\end{center}

Aqua, Ayame, Shion, Korone và Okayu bắt đầu dò tìm nguyên liệu để làm bánh sinh nhật cho Kuon. Cả nhóm bước vào khu thực phẩm, giữa những quầy hàng tràn ngập hương thơm ngọt ngào của bơ sữa và vani.

Nàng Hầu gái, với tinh thần đầy quyết tâm, giơ nắm tay lên cao: ``Được rồi! Trước tiên, ta sẽ tìm nguyên liệu để làm bánh kem (ケーキ) thôi nhỉ!''

Okayu liếc nhìn cô, rồi nhếch mép cười tinh nghịch. Không nói không rằng, nàng Miêu bước tới, nắm lấy vai Ayame và kéo cô nàng Quỷ nhỏ về phía trước mặt ``bạn đời'': ``Bánh hả?''

Nàng Hầu gái chớp mắt vài giây, rồi\dots\ ``\textbf{Cái đó là quỷ (オニ) mà!}''

Nàng Mèo tím bật cười khúc khích, trong khi Ayame giãy giụa: ``Oi oi, sao tự nhiên kéo tôi ra làm gì hả!?''

Không để bị trêu đùa thêm, Aqua quay sang Ayame, quyết định chuyển câu hỏi cho cô nàng. Nhưng không ngờ, Ojou-san nhanh nhảu lấy ra một thứ khác không phải là bánh: ``Bánh sao?''

Aqua nhìn vật trong tay Ayame, rồi xụ mặt: ``Đấy là bánh quy (クッキー) mà! Tôi đang nói đến bánh kem cơ!!''

Nàng Hầu gái thở dài, đưa mắt sang Korone, hy vọng nàng Chó ít nhất sẽ hiểu ý của mình. ``Korone-chan, em có hiểu ý chị nói không đấy?''

Korone không đáp, chỉ cười toe toét rồi lấy ra một cái hình nộm mini. Aqua lập tức trừng mắt, gân xanh nổi lên trên trán: ``Cái đó là chibi (ちーび) mà! Nó còn chẳng phải là thức ăn nữa!''

Cả nàng Miêu và nàng Quỷ sừng cười khúc khích với phản ứng tức giận của nàng Hầu gái. Giống như lần mà họ cùng nhau chơi vần ``cua'' lúc đi bắt cua ở ngoài khơi vào năm ngoái vậy\footnote{Xem lại {\flaregothic ÉPISODE 41}, quyển số Ba}.

Nhưng cuộc tra tấn tinh thần của Aqua vẫn chưa kết thúc.

Shion lặng lẽ tiến đến, khuôn mặt đầy vẻ đắc ý, như thể cô vừa tìm ra một thứ vô cùng quý giá. Nàng Phù thủy lôi ra một viên ngọc hồng sáng lấp lánh, giơ trước mặt Aqua: ``Vậy còn cái này?''

Nàng Hầu gái gần như phát điên, ôm đầu hét lên: ``Cái đó là hồng ngọc (ルビー) mà!! \textbf{Bà kiếm cái đó ở đâu vậy hả!?}''

Shion chỉ tay về một gian hàng đá quý phía xa, cười đầy tự hào: ``Kia kìa.''

Aqua muốn gục ngã tại chỗ: ``Chúng ta đang đi tìm nguyên liệu làm bánh cơ mà\dots''

Korone cười tủm tỉm trước nàng Hầu gái đang gục đầu xuống sàn, tỏ vẻ cảm thông: ``Thôi, không trêu senpai nữa. Nghiêm túc nào.''

\ct

Okayu khoanh tay, nheo mắt nhìn Aqua: ``Nhưng mà, Aqua-san này\dots\ chị có biết làm bánh không?''

``T-Tất nhiên rồi!'' nàng Hầu gái ưỡn ngực tự tin, nhưng rồi ánh mắt lảng đi chỗ khác.

Shion liếc nhìn Aqua một cách đầy hoài nghi: ``Bà chắc chứ?''

``\dots Ờ thì, làm bánh cũng chỉ là nướng bột với đường thôi đúng không\dots?''

Nàng Miêu nghe vậy phì cười: ``Hầy. Được rồi, thôi thì để em với Koro-san tìm nguyên liệu giúp chị vậy.''

Nói rồi, cô kéo tay cô bạn đi về phía quầy bơ sữa, trong khi Ayame và Shion cũng bắt đầu tìm đường đến khu bột mì. Còn Aqua\dots\ cô vẫn đứng đó, mặt đỏ bừng, lẩm bẩm gì đó như thể đang tự trấn an bản thân.

Và thế là, nhóm đầu tiên chính thức bước vào hành trình mua nguyên liệu làm bánh, dù rằng\dots\ chặng đường có vẻ không mấy suôn sẻ.

\ct

Okayu và Korone rảo bước đến khu bơ sữa và hương liệu, dạo quanh những kệ hàng đầy ắp nào là bơ lạt, bơ mặn, vani, sô-cô-la, matcha, dâu tây\dots\ và nhiều hương liệu khác. Cả hai nhìn nhau, vẻ mặt tràn đầy suy tư như đang đứng trước một quyết định hệ trọng trong đời.

``Nè Ogayu, bà nghĩ Kuon-senpai thích vị gì nhất?'' nàng Khuyển hỏi, tay lướt qua một chai chiết xuất hạnh nhân.

``Hmm\dots\ Kuon-senpai à\dots?'' nàng Miêu chống cằm, mắt lim dim như đang chìm vào suy ngẫm. ``Anh ấy thuộc dạng `thực thần', cái gì cũng ăn được, thế nên\dots''

Như thể họ biết được sẽ làm gì tiếp theo, một ý tưởng điên rồ nảy lên trên đầu họ. Cả hai đồng thanh: ``\textbf{Chúng ta mua hết mọi hương liệu luôn đi!}''

Không chút chần chừ, Korone vung tay gom sạch hàng loạt chai lọ trên kệ, Okayu cũng nhanh chóng bê thêm một rổ đầy chai si-rô và tinh dầu thơm các loại.

Những vị khách gần đó nhìn cảnh tượng này đầy ngỡ ngàng, mồ hôi đầm đìa, có người còn suýt đánh rơi cả chai sữa đang cầm trên tay.

``\vie{Ch-chờ đã, hai cô gái đó định làm gì thế kia\dots!?}''

``\vie{Chắc không phải họ định cho tất cả hương liệu vào cùng một cái bánh đâu nhỉ!?}''

``\vie{Họ định mở một tiệm bánh hay sao vậy!?}''

Nhưng cặp đôi Chó-Mèo vẫn bình thản như không, tiếp tục lấy tất cả các hương vị có thể có trong khu này.

``Bà đúng là thông minh đấy, Koro-san\~{}'' nàng Miêu mỉm cười, giọng đầy trìu mến. ``Tôi thật sự cảm động vì bà hiểu rõ tâm lý đàn ông đến vậy\~{}''

``Ehehe\~{}Đương nhiên rồi!'' nàng Khuyển ưỡn ngực tự hào, rồi bất ngờ nháy mắt. ``Nhưng mà bà cũng đáng yêu ghê đó, Ogayu\~{}''

``Hể\~{} Bà tính tán tỉnh tôi ngay trong siêu thị luôn à?''

``Ơ, chẳng lẽ bà không thích sao?''

Okayu nhẹ nhàng tựa vào Korone, hơi thở phả nhẹ bên tai: ``Ai mà biết được nhỉ\~{}''

Ngay lập tức, khách hàng xung quanh đồng loạt rùng mình. Một ông chú đứng gần đó vội kéo áo vợ mình, thì thầm: ``\vie{Mình ơi, bọn trẻ con bây giờ đáng sợ quá\dots}''

Còn một chị gái trẻ tuổi thì quay đi chỗ khác, che mặt: ``\vie{Tôi không chịu nổi nữa\dots\ Đây là siêu thị cơ mà, không phải quán cà phê hẹn hò đâu\dots}''

Nhưng Okayu và Korone? Cả hai vẫn vô tư gom thêm đồ, tiếp tục bàn bạc xem có nên bỏ thêm\dots\ hương vị sầu riêng vào bánh không. Và thế là, những người mua sắm gần đó âm thầm di chuyển ra xa, tránh để bị cuốn vào cái vòng xoáy ``năng lượng kỳ lạ'' của hai cô gái \textit{GAMERS} này.

\ct

Cặp đôi tiếp tục tiến về quầy bơ sữa, nhưng lần này, thay vì chỉ lấy một hai hộp, họ\dots\ quơ sạch cả một ngăn tủ.

``Này Ogayu, chúng ta cần bao nhiêu bơ nhỉ?'' Korone cầm một hộp bơ lạt, nhưng rồi lại lấy thêm ba hộp nữa.

``Hmm\dots\ Bà còn phải hỏi à? Cứ lấy nhiều nhiều đi, lỡ thiếu thì sao?''

Nói rồi, Okayu chộp luôn cả chục hộp bơ, quẳng hết vào giỏ hàng.

``Sữa thì sao?''

``Ờm\dots\ Chúng ta định làm bánh bông lan kem đúng không-? Vậy thì ít nhất cũng phải mua\dots''

Okayu chớp mắt một cái, rồi vơ nguyên ba hộp sữa loại lớn.

``\dots Chắc vẫn chưa đủ đâu nhỉ?''

Korone nhìn chỗ sữa trong giỏ, rồi lại nhìn Okayu, sau đó không nói không rằng\dots\ bê hẳn một thùng nguyên liệu chuyên dụng dành cho tiệm bánh.

``Có khi tôi nên mở tiệm bánh luôn ấy chứ!'' nàng Miêu bật cười, vỗ nhẹ lên hộp bơ trên tay.

``Tôi sẽ là người thử bánh! Được làm bởi Ogayu đáng yêu của tôi thì chắc chắn sẽ ngon lắm!''

``Hể\~{} Bà khen tôi như thế, có khi tôi sẽ cố ý làm bánh hỏng để bà ăn đó nha\~{}''

``Hỏng cũng ăn! Tình yêu của tôi dành cho bánh của Ogayu là vô điều kiện mà!''

Câu đối đáp của hai người khiến một nhân viên siêu thị đi ngang qua lặng lẽ thở dài, tự nhủ: ``\vie{\textit{Hai vị khách này có thể đừng thả bầu không khí yêu đương ngọt lịm ngay trong siêu thị được không vậy\dots?}}''

Sau một hồi bàn bạc, cả hai quyết định:

``Chúng ta sẽ làm ba tầng!''

``Ừ! Lớp đầu tiên sẽ là vani, lớp thứ hai sô-cô-la, lớp trên cùng dâu tây!''

``Sau đó phủ kem lên, rồi rắc thêm\dots\ hmm\dots\ mấy thứ hương liệu chúng ta vừa mua!?''

``Chính xác! Tạo ra một chiếc bánh độc nhất vô nhị!''

``Kuon-senpai mà biết được như vậy sẽ bất ngờ lắm đấy!''

Và thế là, Okayu và Korone gật gù, vô tư gom thêm nguyên liệu, hoàn toàn không để ý đến những ánh mắt khó hiểu của những khách hàng xung quanh.

\ct

``Nào, xong rồi!'' nàng Khuyển vác trên vai một túi đầy chai lọ, còn cô bạn cũng xách hai giỏ nặng trịch.

``Ừm, chắc thế là đủ\dots\ Tôi nghĩ với chỗ này thì đến Kuon-senpai cũng phải bối rối rồi.''

``Ha ha ha\~{} Tôi thật sự mong chờ lúc cậu ta thử miếng bánh đầu tiên!''

Nói xong, cả hai bật cười đầy nham hiểm, rồi thong thả tiến về quầy thanh toán. Và thế là, quầy bơ sữa của siêu thị đã chính thức bị ``càn quét'' sạch sẽ chỉ trong vòng chưa đầy mười phút!

\begin{center}
    \uline{Quầy bột mì.}
\end{center}

Ở một diễn biến khác\dots

Ayame và Shion đứng trước gian hàng bột mì, khoanh tay suy nghĩ như thể đang bàn chuyện quốc sự.

``Được rồi, để làm bánh chúng ta cần gì nữa nhỉ?'' nàng Phù thủy gõ gõ cằm, liếc nhìn giá nguyên liệu như thể đang tìm kiếm một món bảo vật cổ xưa.

Nàng Quỷ sừng khoanh tay tràm ngầm: ``Hmm\dots\ Xem nào\dots\ Bột mì này, men nở này\dots''

``Ừ, đường nữa, muối một tí, hương liệu với bơ sữa chắc Okayu với Korone lo rồi\dots''

Shion gật gù, rồi đột nhiên\dots\ chìa tay ra trước mặt Ayame. ``Bà làm cái gì vậy?''

``Xin một giọt máu quỷ để hoàn thiện công thức phép thuật làm bánh siêu cấp.''

Ayame phồng má lên, tỏ vẻ bất mãn: ``Bà nghĩ tôi là ai hả? Tôi là Quỷ, không phải là cái bình máu di động của bà đâu nha!''

Shion cười phá lên: ``Ồn quá\~{} Đùa tí thôi mà\~{}''

Khách hàng gần đó đã bắt đầu nhìn hai người bằng ánh mắt khó hiểu, nhưng hai nàng vẫn tiếp tục bàn bạc.

\ct

Họ đi tới một quầy bột mì, với đủ kiểu loại khác nhau và đủ loại nhãn hiệu khác nhau. Trông không khác gì một ma trận bột mì cho khách hàng lựa chọn cả.

Ayame lấy ra một gói bột mì, rồi cẩn thận nhìn từ trên xuống dưới: ``Này, bột mì này có gì khác bột mì kia không?''

``Hmm\dots\ để tôi xem\dots'' Shion liếc sang một gói bột mì khác: ``Ôi chà, giá cái này rẻ hơn này!''

``Ể, rẻ hơn!? Vậy lấy cái này đi!''

``Khoan! Bà phải xem nhãn hiệu chứ!''

``Ừ nhỉ\dots\ Ủa khoan, bà nói vậy tức là bột mì cũng có cấp bậc khác nhau à!?''

``Ừ, giống như phép thuật ấy. Có loại chỉ để làm bánh mì, có loại để làm bánh bông lan.''

``Ể\dots\ tức là nếu tôi lấy cái này\dots'' Nói rồi Ayame vớ bừa một bịch bột trên kệ. ``\dots Là tôi có thể làm phép triệu hồi bánh ra từ hư vô đúng không!?''

Nàng Phù thủy thở dài, đổ mồ hôi hột: ``Không phải kiểu đó!''

Những người xung quanh chỉ thấy hai cô gái cứ đứng đó, chỉ trỏ vào mấy bịch bột mì rồi lẩm bẩm những câu thoại nghe không khác gì đang thảo luận ma thuật hắc ám. Một nhân viên siêu thị đi ngang qua, nghe lỏm được vài câu liền nhanh chóng bỏ đi, không muốn bị kéo vào cái cuộc hội thoại ``quá siêu hình'' này.

Cuối cùng, sau một hồi tranh luận sôi nổi về ``bột mì cao cấp'' và ``bột mì bình dân'' khiến người ngoài nghe còn muốn ngất xỉu, Ayame và Shion cũng gom đủ nguyên liệu cần thiết, trên tay mỗi người cầm ít nhất ba bốn túi bột lớn.

``Được rồi! Giờ thì tiếp theo là gì nhỉ?''

``Hmm\dots\ Chúng ta có cần một quyển sách phép thuật hướng dẫn làm bánh không ta?''

``Bà khỏi lo, tôi có công thức trên điện thoại rồi!''

``Hể\~{} Nhưng dùng ma thuật làm bánh chắc nhanh hơn nhỉ\~{}''

Nàng Quỷ sừng nhe răng cười, châm chọc cô bạn: ``Lần cuối bà dùng phép thuật là lần nào ấy nhỉ? À đúng rồi, lần suýt làm nổ cái bếp của nhà bà ấy!''

Tới đây, Shion đỏ mặt, rồi rít lên như không muốn nhớ lại khoảnh khắc đáng xấu hổ ấy: ``Im đi!''

\ct

Sau khi chất đầy bột mì trong giỏ hàng, Ayame và Shion chuyển sang gian hàng phụ gia. Lần này, họ tiếp tục cuộc ``nghiên cứu'' một cách cực kỳ nghiêm túc, nhưng theo một phong cách\dots\ đáng ngờ hơn.

``Được rồi, tiếp theo là đường này. Bà thích loại nào?'' nàng Phù thủy hỏi, tay lướt dọc theo các kệ hàng.

``Để Ojou xem nào\dots''

Nàng Quỷ sừng cầm một túi đường trắng lên, chăm chú nhìn nó dưới ánh đèn siêu thị. Rồi đột nhiên\dots ``Bà nghĩ đường này có phải là hàng thật không?''

``Hể? Bà nói gì vậy?'' Shion ngạc nhiên.

``Bà không thấy mấy viên đường này sáng bóng một cách đáng ngờ à? Nó có khi nào là giả không?'' Ayame trố mắt, chỉ trỏ vào những hạt đường bên trong gói đường.

Shion nheo mắt nhìn túi đường trong tay Ayame, rồi bất chợt run lẩy bẩy: ``Giờ bà làm tôi hoang mang theo rồi này!''

``Tôi nghe nói có những nơi bán đường giả đấy\dots''

``Hả!?''

``Nếu chúng ta vô tình mua nhầm thì sao?''

Hai người cùng quay ra nhìn kệ hàng một cách đầy nghi ngờ, như thể đang truy tìm một âm mưu bí ẩn nào đó. Một vị khách gần đó đứng hình mất mấy giây, rõ ràng không hiểu hai cô gái này đang nói cái gì.

``Thế còn loại này?'' nàng Phù thủy cầm một túi đường nâu lên.

``Đường nâu à\dots\ Có khi nào đây là phiên bản tăng cấp của đường trắng không?''

``Bà nghĩ đây là game nhập vai hả!?''

Sau một hồi bàn luận như thể đang điều tra một vụ án, Ayame và Shion cuối cùng cũng chọn bừa một túi đường và lấy thêm một túi nữa, nhưng vẫn không ngừng lẩm bẩm về độ đáng tin của nó.

Tiếp theo, họ tới khu vực muối. Nhưng lần này, vấn đề lại xoay quanh lượng muối cần dùng.

``Tôi nghĩ ta nên mua một ít muối,'' Shion liếc nhìn.

``Ừm, nhưng mà ít là bao nhiêu?''

``Chắc chỉ cần một nhúm thôi?''

Ayame suy nghĩ một lúc rồi\dots\ với tay lấy hẳn một bao muối nặng một cân.

``Một nhúm cỡ này chắc đủ nhỉ, yo?''

``Bà định làm bánh hay dựng lò luyện đan vậy!?''

Một khách hàng gần đó đang định lấy một gói muối, nhưng sau khi nghe cuộc đối thoại của hai người, ông ta lặng lẽ rút tay lại và rời đi trong im lặng.

Tiếp đến là tinh dầu vani. Nhưng khi tới kệ hàng, họ lại vướng vào một cuộc tranh luận mới.

``Ủa, sao có nhiều loại thế?'' Shion khẽ choáng ngợp với số lượng loại tinh dầu nhìn thấy trên kệ.

``Để tôi xem nào\dots\ À-rế?''

Nàng Phù thủy nghiên đầu hỏi: ``Sao thế Ayame?''

Nàng Quỷ sừng cầm một lọ tinh dầu cỡ lớn lên, tò mò đáp: ``Bà nhìn nè, cái này có ghi là `hương liệu tổng hợp' này!''

``Hể!? Tức là đây là hương vani giả hả!?''

``Vậy còn cái này?'' Ayame cầm một lọ khác lên.

``Chiết xuất vani nguyên chất\dots\ Nghe như đồ cao cấp nhỉ?''

``Bà nghĩ nó có phải hàng thật không?'' nàng Quỷ sừng nhíu mày, nhìn lọ tinh dầu với vẻ nghi ngờ.

``Tất nhiên là hàng thật rồi! Đừng có nghi ngờ mọi thứ như thế chứ!'' Shion rít lên, bắt đầu mất kiên nhẫn.

Một lần nữa, khách hàng xung quanh chỉ biết nhìn hai cô gái vừa bàn bạc vừa trừng trừng nhìn chằm chằm vào từng món hàng như thể chúng là cổ vật quý hiếm. Cuối cùng, sau khi trải qua một chuỗi suy luận đầy hoang mang, họ chọn đại một chai tinh dầu vani và nhanh chóng chất hết mọi thứ vào giỏ.

Họ cũng mua thêm vài nguyên liệu làm bánh nữa như trứng (họ mua hai tá), bột nở (hai gói lớn), và vài nguyên liệu lỉnh kỉnh khác.

Và thế là hai cô phù thủy và quỷ nhỏ tiếp tục lảm nhảm trên đường đến quầy tính tiền, để lại đằng sau những vị khách vẫn còn đang cố gắng giải mã xem rốt cuộc hai người kia đang làm cái trò gì ở siêu thị này.

\ct

\begin{center}
    \uline{Quầy sách.}
\end{center}

Aqua lủi thủi đi dọc theo gian hàng sách, ánh mắt lấp lánh khi nhìn vào từng cuốn sách nấu ăn bày trên kệ.

``\textit{Mình nhất định phải mua một cuốn! Hừm, không thể để bọn họ coi thường mình được!}'' cô lẩm bẩm một mình, nhớ lại cảnh Okayu nhìn mình với ánh mắt đầy hoài nghi về trình độ nấu ăn. Không thể chịu được điều đó, nàng hầu gái quyết tâm ``nâng cấp'' kỹ năng bếp núc - một trong những kỹ năng VỐN CÓ của một hầu gái.

Sau một hồi lục lọi, Aqua vớ lấy một cuốn sách có tựa đề ``Hướng dẫn làm bánh cơ bản dành cho người mới bắt đầu''.

Cầm được cuốn sách trên tay, cô hí hửng thầm nghĩ: ``\textit{Được rồi! Mình sẽ học từ đầu! Từ giờ, Aqua-sama sẽ là bậc thầy làm bánh!}''

Vừa lúc đó, một nữ nhân viên siêu thị thân thiện bước tới, nhìn Aqua cười hiền: ``\vie{Ô, em đang tìm sách nấu ăn à? Cần chị tư vấn không?}''

Aqua giật mình, ú ớ trước nữ nhân viên đó: ``À, dạ, à thì là\dots\ Em\dots''

Nữ nhân viên vẫn giữ cho mình một nụ cười thân thiện: ``À, em nói tiếng Nhật à? Không sao đâu, chị hiểu.''

Thấy thế, nàng Hầu gái thở phào nhẹ nhõm. Người nữ nhân viên nhắc lại câu hỏi: ``Em tới đây tìm sách nấu ăn sao? Cần chị tư ván không?''

Aqua gật đầu lia lịa: ``V-vâng! Em muốn học cách làm bánh sinh nhật!''

``Ồ, dễ thôi! Nhưng mà\dots\ em định chỉ làm bánh thôi sao?''

Nữ nhân viên nghiêng đầu, rồi bất ngờ nở một nụ cười đầy ẩn ý: ``Hay là em cũng muốn học cách nấu các món ăn ngon khác cho bạn trai nữa?''

``Ể!?'' mặt của nàng Hầu gái lập tức chuyển sang màu cà chua chín.

Cô nàng khươ tay múa chân, lắc đầu nguầy nguậy: ``K-không phải! Không phải như chị nghĩ đâu!''

``Ôi chà, không cần ngại mà! Nếu em muốn gây ấn tượng với người ấy, chị khuyên em nên học thêm mấy món như omurice, cơm cà ri, hay là mấy món ăn gia đình ấy!''

``Không phải! Không phải thế đâu, chị ơi!'' nàng Hầu gái quơ tay loạn xạ, cố gắng đính chính nhưng càng nói càng rối.

``Thật đấy! Đây là\dots\ cho một người bạn! Không phải như chị nghĩ -đâu mà!''

``Ara\~{} Một người bạn quan trọng đến mức em phải học nấu ăn cho cậu ấy à?''

Tới đây, Aqua hét lên trong lành, đầu bắt đầu xì khói như ấm nước đun sôi vì ngượng ngùng. Cô nàng bối rối tới mức suýt nữa đánh rơi cả cuốn sách trên tay. Cô vội vã vớ lấy sách, cúi đầu cảm ơn rồi chạy trối chết ra khỏi khu vực đó, để lại nữ nhân viên siêu thị cười thầm đầy thích thú.

``\vie{Thanh xuân thật là đáng yêu thật á\~{}!}''

Aqua lao ra khỏi quầy sách với tốc độ của một cơn lốc, hai má vẫn còn đỏ bừng như thể vừa bị ai đó trêu chọc quá đà.

``\textit{M-mình mà lại nấu ăn cho Kuon-senpai á!? Không đời nào! Không bao giờ!}

Nhưng ngay khi vừa nói ra, trong đầu cô bất giác hiện lên một hình ảnh đầy mơ hồ: cô trong bộ tạp dề, đứng trong bếp, cẩn thận múc từng thìa canh nóng hổi ra bát, rồi nhẹ nhàng đặt xuống trước mặt cậu Tổng.

``Đây, anh ăn thử đi\dots\ Em đã nấu món này cho anh đó\dots''

``Được thôi\dots\ Không biết món ăn mà Akutan làm sẽ thế nào nhỉ\dots''

Aqua giật bắn người, hai tay ôm đầu, quằn quại ngay giữa lối đi của siêu thị như thể vừa bị ai đó niệm chú.

Và rồi, cô hét lớn trong vô thức: ``\textbf{Không đời nào có chuyện đó xảy ra đâu mà!!!}''

Hiển nhiên, tiếng hét của cô khiến cho toàn bộ khách hàng xung quanh quay đầu nhìn cô với ánh mắt khiến họ ``lú như con cúc cu''. Một ông bác đang chọn rau củ đứng đơ ra, một bà cụ khẽ thở dài, còn một đứa trẻ nhỏ kéo áo mẹ thì thầm: ``\vie{Mẹ ơi, chị kia bị sao thế ạ\dots?}''

Nhận ra mình vừa làm điều gì đó quá lố, Aqua chết lặng trong hai giây, rồi nhanh như chớp, cúi đầu thật sâu và hét lên: ``\textbf{C-Cháu xin lỗi!!!}''

Ngay sau đó, cô lập tức co giò bỏ chạy về phía quầy tính tiền, mặt đỏ như gấc chín.

\dots Nhưng, trong thâm tâm, có lẽ cô cũng muốn làm như vậy thật. Một ngày nào đó, cô sẽ gây ấn tượng cho cậu Tổng, người mà cô quý trọng nhất.

\vspace{-40pt}
\subsubsection*{\phantomsection\label{P1.1b}}
\addcontentsline{toc}{subsubsection}{Đội \#{\digi 86432E} - Lamy, Roboco và Haato}

\begin{center}
    {\color{T2} Nhóm 2 - Lamy, Roboco và Haato\\Đồ ăn vặt và đồ uống.}
\end{center}

Lamy, Roboco và Haato bước vào khu vực đồ uống, mắt nhìn khắp các kệ hàng đầy những chai lọ rực rỡ sắc màu.

Nhưng với cô nàng Yêu tinh tuyết bợm rượu, thứ mà cô để ý nhiều nhất lại là\dots\ những thức uống có cồn.

``Để em lấy cái này, cái này\dots\ cả cái này nữa!'' cô nàng hào hứng với tay lấy hết chai này đến chai khác, hai mắt sáng rỡ như một đứa trẻ lạc vào cửa hàng kẹo.

Haato nhìn đống chai rượu dần chất đầy trên xe đẩy mà mặt tái mét. Cô hoảng hồn vội căn ngăn: ``T-Từ từ đã nào, Lamy-chan! Sao em mua lắm rượu thế hả!?''

``Ở đây có nhiều loại rượu mà em chưa thử bao giờ, em muốn uống thử ghê!'' Lamy đáp, đôi mắt vẫn sáng như sao, không thèm để tâm tới tiền bối cô.

``Cái đó không quan trọng! Em biết là chúng ta đang mua đồ uống cho sinh nhật của Kuon-senpai chứ!?''

Roboco cũng gật đầu phụ họa: ``Em ấy nói đúng đó, Lamy-chan à. Chúng ta nên mua cho anh ấy trước đi đã chứ.''

``Đúng đó\dots'' nàng Song tính thở dài, rồi liếc về giọng nói vừa cất lên kia.

``Haato-chan nhìn chị chằm chằm thế?'' nàng Người máy nghiêng đầu khó hiểu.

``Tay chị\dots\ Tay chị cầm hộp cà ri là sao kia!?'' nàng Song tính thảng thốt.

Roboco nhìn xuống tay mình, nơi đang giữ một hộp cà ri ăn liền. Cô nhún vai thản nhiên: ``Ờ thì, chị thấy ngon mà.''

Haato trừng mắt: ``Cái\dots\ Cà ri đâu phải đồ uống!?''

``Nhưng nó là súp cà ri mà?'' Roboco ngây thơ đáp lại.

``\textbf{Đó không phải là vấn đề!!}'' nàng Trung học gắt lên, rồi thở dài chán chừng. Cô cảm thấy cái nhóm này đúng là chẳng có tí trách nhiệm nào cả.

Tới lúc này, nàng Yêu tinh tuyết mới thoát khỏi cơn mơ rượu của bản thân, lập tức hỏi: ``Cơ mà\dots\ Kuon-senpai thích uống gì nhỉ?''

``Chị nghĩ là gì cũng được,'' Roboco nhớ lại. ``Sora-chan với Kuon-senpai có nói thế mà.''

``Hồi Tết năm ngoái, em có nhớ Kuon-senpai ưa uống nước ép hoa quả để giữ dáng mà nhỉ?'' Haato bổ sung.

Thấy vậy, Lamy vỗ tay cái bép, như thể nhận ra chân lý. ``Vậy là rượu trái cây được đúng không!?''

``\textbf{Không phải như thế!!}'' nàng Song tính gào lên, bắt đầu mất kiên nhẫn với cô nàng ham rượu này.

Cô ôm đầu bất lực. Nhưng ngay khi cô đang định giải thích lại, Lamy chợt cầm một cái chai kỳ lạ lên.

Trên nhãn chai, dòng chữ lớn in đậm: ``Gì cũng được.''

Nàng Yêu tinh tuyết nhìn nó chằm chằm, rồi hỏi một cách đầy nghiêm túc: ``Có phải chai này không?''

Roboco nhìn chai rượu, bật cười khúc khích: ``Chị nghĩ là nó đó!''

``\textbf{Không phải thế!!}'' tiếng thét tuyệt vọng của Haato vang lên. Những khách hàng gần đó giật mình quay lại nhìn ba cô gái kia mà đổ mồ hôi hột.

\ct

Sau màn náo loạn với chai rượu ``Gì cũng được'', cuối cùng nhóm cũng bắt đầu nghiêm túc chọn đồ uống.

``Lần này, rút kinh nghiệm từ vụ Tết năm ngoái, chúng ta sẽ không mua đủ thứ tạp nham nữa,'' Haato khoanh tay, nghiêm túc lên tiếng. ``Chúng ta chỉ mua hai, ba loại nước theo sở thích chung của mọi người, với số lượng vừa đủ thôi!''

``Được rồi được rồi, hiểu rồi,'' Roboco gật gù. ``Vậy thì\dots\ Một loại nước ngọt, một loại nước ép và nước lọc là ổn nhỉ?''

Nàng Song tính búng tay, cười lớn: ``Ý hay đó!''

Lamy cũng theo đà đó mà lên tiếng vói nụ cười trên mổi: ``Và cả rượu nữa chứ!''

``Tất nhiên là không được rồi!'' Haato quay ngoắt về phía cô, nói với giọng nghiêm túc.

Nàng Yêu tinh tuyết đứng hình: ``\dots Ể?''

``Chúng ta đã thống nhất là không được uống rượu bừa bãi mà, đúng chưa?'' Haato nhắc lại, nghiêm giọng như một người mẹ quyết không chiều hư con cái mình.

``N-Nhưng mà\dots''

Chưa để Haato kịp phản ứng, cô nàng Yêu tinh tuyết đổ gục xuống sàn siêu thị, nằm lăn lóc ăn vạ như một đứa trẻ. Cô vừa lăn lóc vừa la khóc: ``\textbf{Không có rượu thì em sống sao đây!!!}''

Roboco nhìn cảnh tượng đó cũng chỉ biết cười trừ, còn Haato thì nhất quyết không chịu thua: ``Dậy ngay!''

``Không có rượu thì làm sao em thể hiện được bản sắc của mình chứ\dots!''

Nàng Trung học bắt đầu mất bình tĩnh, lớn giọng: ``\textbf{Dậy ngay, Lamy!!}''

Nhưng Lamy dường như không có dấu hiệu gì muốn chịu thua cả: ``Không có rượu thì em không còn là Yukihana Lamy nữa!''

``\textbf{Đừng có mà trẻ con nữa!}'' nàng Song tính càng la lớn hơn.

Khách hàng xung quanh dừng lại nhìn nàng Yêu tinh xoay người ăn vạ, rồi lăn lóc khắp sàn nhà với những dấu hỏi đặt trên đầu. Có người lẩm bẩm:

``\vie{Con gái thời nay đáng sợ thật\dots}''

``\vie{Mình có nên gọi nhân viên siêu thị không nhỉ?}''

``\vie{Ủa, hình như cô kia là thần tượng đấy?}''

Haato ôm đầu, cảm thấy cả linh hồn lẫn thể xác đều kiệt quệ, hoàn toàn bất lực với cô nàng hậu bối này. Cuối cùng, cô đành phải xuống nước: ``Được rồi, được rồi! Cho em mua rượu, được chưa!?''

Lamy bật dậy như chưa từng có gì xảy ra. Cô reo hò vui sướng, giơ hai tay lên trên cao: ``Yatta!''

Roboco vỗ vai nàng Song tính đầy cảm thông: ``Đấu tranh với con nghiện rượu chỉ tổ tốn sức thôi. Em không lại với em ấy được đâu.''

Haato không còn sức mà đáp lại tiền bối Người máy được nữa.

Và thế là nhóm chính thức chốt đơn: một loại nước ngọt, một loại nước ép, một ít nước lọc, và đặc biệt\dots\ một chai rượu dành riêng cho Lamy.

\ct

Sau khi đã ``chốt đơn'' phần đồ uống, Haato và Roboco bắt đầu bàn bạc về các món ăn vặt cho bữa tiệc tối. Trong khi đó, Lamy đang đứng ở một góc, ôm khư khư chai rượu như thể đó là một loại quốc bảo quý hiếm.

``Nào, giờ đến đồ ăn vặt!'' nàng Người máy vỗ tay. ``Sau bữa chính thì chắc chắn ai cũng muốn có gì đó nhâm nhi, đúng không?''

``Chuẩn luôn!'' Haato hăng hái. ``Mình sẽ cần cả đồ ngọt và đồ mặn. Đồ ngọt thì bánh kẹo là chắc chắn, nhưng còn đồ mặn thì sao nhỉ?''

``Bim bim?''

``Quá bình thường!''

``Khoai tây chiên?''

``\textbf{Quá} bình thường!''

``Vậy\dots\ chân gà cay?''

``\textbf{Quá-}''

``Thôi đừng la nữa, mọi người đang nhìn kìa.''

Quả nhiên, một số khách hàng gần đó đã bắt đầu liếc về phía hai cô gái với ánh mắt pha lẫn tò mò và lo ngại.

``Thế em muốn ăn gì?'' Roboco khoanh tay.

Nàng Song tính đặt tay lên cằm, ra vẻ suy tư cực kỳ nghiêm túc: ``\dots Thịt bò sấy.''

``Hả!?''

``Một bữa tiệc sinh nhật hoành tráng thì không thể thiếu thịt bò sấy được!''

``Nhưng mà\dots\ bình thường ai lại ăn cái đó sau bữa tối chứ?''

``Em á!''

``\dots'' nàng Người máy đơ mắt vài giây, bị choáng với lời tự thú của đàn em.

Cô cũng chẳng buồn không đào sâu vào tâm lý phức tạp của Haachama nữa. Có lẽ cuộc cãi cọ giữa cô ấy và Lamy ban nãy đã làm Haato không còn tỉnh táo.

``Được rồi, cứ lấy một ít vậy.''

``Được quá-chama!''

``Mà khoan\dots\ Lamy-chan, em nghĩ sao?''

Hai cô nàng quay lại thì thấy nàng Yêu tin tuyết vẫn đang âu yếm chai rượu, thì thầm những lời yêu thương như đang nói chuyện với một người yêu dấu: ``Chúng ta cuối cùng cũng đến với nhau\dots\ Định mệnh thật kỳ diệu\dots''

Roboco và Haato im bặt, nhìn đàn em của họ ôm ấp món bảo bối của cô mà không nói nổi một câu nào.

Nàng Song tính búng tay trước mặt Lamy, nwhng không nhận lại phản ứng nào.

``\dots Hết cứu rồi,'' cô tặc lưỡi.

``Em ấy có vẻ đang ở một chiều không gian khác rồi,'' nàng Người máy cười trừ.

``Kệ đi, giờ em muốn tìm thêm bánh kẹo!''

Roboco nhún vai, rồi cùng Haato tiếp tục lùng sục khu đồ ăn vặt. Sau một hồi tranh cãi dữ dội giữa bánh bông lan với bánh quy (Haato thắng), họ cũng gom đủ chiến lợi phẩm: một ít bánh ngọt, một ít thịt bò sấy, và vài món ăn vặt khác để nhâm nhi.

Lamy, cuối cùng cũng chịu rời khỏi thế giới riêng của mình, quay lại với nhóm cùng ánh mắt lấp lánh: ``Hai người mua gì đó?''

``Đồ ăn vặt!'' Roboco đáp gọn.

``Có món nào hợp với rượu không?''

``Có chứ!'' nàng Người máy nói, cầm những hộp bánh ăn vặt trước mặt nàng Yêu tinh tuyết.

Lamy đặt ngón tay lên môi, mỉm cười đầy bí ẩn: ``Fufufu\dots\ Tất nhiên rồi\~{}''

Haato và Roboco chỉ còn biết lắc đầu, trong khi Lamy tỏ ra vô cùng mãn nguyện.

Và thế là nhóm hoàn tất phần nhiệm vụ của mình, sẵn sàng tiến về quầy tính tiền.

\vspace{-40pt}
\subsubsection*{\phantomsection\label{P1.1c}}
\addcontentsline{toc}{subsubsection}{Đội \#{\digi 923E7\ 1} - Noel, Pekora, Rushia, Flare }

\begin{center}
    {\color{T3} Nhóm 3 - Noel, Pekora, Rushia và Flare \\Thịt}
\end{center}

Nhóm Noel, Pekora, Rushia và Flare tiếp tục nhiệm vụ của mình: săn thịt! Trong siêu thị.

Họ đứng trước quầy thịt, nơi đủ loại thịt bò, thịt lợn, thịt gà được bày biện đẹp mắt dưới ánh đèn sáng rực. Hương thơm của thịt tươi khiến Hiệp sĩ gần như phát cuồng.

``Không biết ta sẽ ăn lẩu hay ăn nướng nhỉ?'' nàng Pháp sư nghiêng đầu, ngắm nghía những bọc thịt được bày biện gọn gàng trên các gian hàng.

``Lẩu cũng có thể đó, nhưng mà trời nóng thế này thì hơi khó chịu, peko\dots'' Pekora vuốt cằm suy nghĩ.

``Tôi thấy ăn lẩu hay nướng đều ổn cả,'' Flare hào sảng đáp. ``Cái ngày mà tôi xuống nhà Kuon-senpai với Mio-senpai và Lamy-chan để học hỏi văn hóa uống rượu\footnote{Xem lại {\flaregothic ÉPISODE 92}, quyển số Bảy.}, thấy phòng khách có điều hòa mà.''

``Tôi thì gì cũng được, miễn là có thịt bò!'' Noel siết chặt nắm tay, mắt sáng rực như một chiến binh sắp ra trận.

Rushia toát mồ hôi hột: ``Bà chỉ quan tâm đến thịt bò thôi à-nanodesu\dots''

Nàng Dị giới lắc đầu cười: ``Mà thịt thôi thì chưa đủ, còn phải có rau nữa chứ.''

Nàng Pháp sư liếc nhìn xung quanh: ``Hình như quầy rau ngay gần đây đúng không nhỉ\dots?''

Ngay lập tức, cả ba người quay sang nhìn Pekora với ánh mắt ``mất nhân tính''.

``C-Chuyện gì thế, peko!?'' nàng Thỏ thảng thốt.

Nàng Hiệp sĩ nở nụ cười đầy ẩn ý, tay vươn tới một củ cà rốt ngay trên tóc Pekora: ``Chúng ta có lấy luôn mấy củ cà rốt này không nhỉ?''

``Ý hay đấy-nanodesu!'' nàng Pháp sư hưởng ứng ngay lập tức.

Ngay lập tức, nàng Thỏ hoảng hốt ôm lấy củ cà rốt gắn trên lọn tóc, như thể đó là của quý: ``Khoan, khoan, \textbf{chúng là của tôi mà, peko!!!}''

``Thôi nào mọi người,'' nàng Dị giới cười hiền, đi tới bảo vệ Trưởng nhóm của Thế hệ thứ Ba, ``đừng trêu cậu ấy nữa. Bọn tôi đùa ấy mà.''

Pekora vẫn ôm chặt củ cà rốt, đôi mắt long lanh như sắp khóc, trong khi ba người còn lại cười sảng khoái.

``Thôi được rồi, chúng ta bàn tiếp vụ thịt đi,'' nàng Yêu tinh vàng nói, trước khi nàng Thỏ thực sự nổi giận.

``Tôi nghĩ chắc đội kia lo vụ rau củ rồi-nanodesu,'' Rushia tiếp lời, ``nên tôi nghĩ ta cứ tập trung vào thịt thôi.''

``Được thôi!'' nàng Hiệp sĩ cười khoái chí. Cô siết chặt lòng bàn tay, rồi giơ lên trời đầy năng lượng: ``Giờ thì\dots\ \textbf{Mua thịt nào!!}''

Không ai có thể cản nổi một ``Hiệp sĩ cuồng thịt'' khi đã vào trạng thái săn mồi.

\ct

Noel đứng trước quầy thịt với ánh mắt rực lửa, tựa như một chiến binh sắp bước vào trận chiến. Cô nắm chặt giỏ hàng, đảo mắt qua từng loại thịt bò được trưng bày đẹp mắt trong tủ kính mát lạnh.

``\textbf{Bắt đầu nào!}'' Noel hét lên, khiến những khách hàng xung quanh cô giật mình đánh rơi cả giỏ hàng.

Không chần chừ, cô nàng Hiệp sĩ cuồng thịt lao tới, quơ lấy những phần thịt bò ngon nhất.

Từ ba chỉ bò, đùi bò, sườn bò, cho tới cả thăn nội và thăn ngoại của thịt bò nữa.

Mỗi lần hô tên một loại, Noel lại chộp lấy vài gói thịt và quăng vào giỏ hàng, tốc độ nhanh đến mức Flare còn chưa kịp nhìn rõ.

``C-Chậm lại đã, bà tính mua hết quầy thịt à!?'' Rushia toát mồ hôi hột trước tốc độ siêu thanh của nàng Hiệp sĩ.

``Dĩ nhiên rồi!'' cô đáp lại đầy kiên quyết. ``Với mình, thịt bò là chân ái!''

``Chắc cậu ta cũng sẽ mua cả lưỡi bò và nội tạng bò nếu họ có bán nhỉ,'' Flare bật cười, nhưng cũng không giấu nổi lo lắng trước cô bạn cuồng thịt bò này.

Noel vẫn chưa có dấu hiệu dừng lại. Thậm chí, khi nhìn thấy một bảng giá đặc biệt, cô vội chạy tới và\dots

``\textbf{Tuyệt quá!! Ở đây còn bán nguyên con bò này!!}''

Cả nhóm chết lặng. Trước mặt họ là một con bò còn sống, đứng nhởn nhơ ngay cạnh quầy thịt với cái bảng đề giá treo ngay cổ nó. Những vị khách đi qua bảng giá đó mà cũng ngỡ ngàng không kém gì, nhất là với một loài vật vẫn còn có thể cất tiếng ``ụm bò~{}'' ngay giữa một nơi có vẻ\dots\ không phù hợp này.

``T-Tại sao siêu thị bán cả con bò này vậy!?'' Rushia ôm đầu hét lớn.

``Tôi nghi là siêu thị này đang có chương trình đặc biệt đó\dots'' Flare vuốt cằm suy nghĩ nghiêm túc. ``Cơ mà bán nguyên con bò còn sống thế này thì tôi chưa thấy bao giờ\dots''

``\textbf{Không! Không đời nào!!}'' Pekora hoảng loạn, khiếp sợ với thứ được rao bán trong một nơi mà không ai ngờ tới.

Noel vỗ tay đầy phấn khích, hớn hở nói với các thành viên trong nhóm: ``Mọi người nghĩ sao!? Nếu mua nguyên con, chúng ta sẽ có đủ thịt cho cả tháng luôn!''

Khuôn mặt của nàng Thỏ và nàng Pháp sư xanh lét với những thứ vô lý trước mắt.

Nhưng, đó vẫn chưa là gì so với sự điềm tĩnh của cô nàng Dị giới - người đáng ra phải là hoảng hồn nhất trong nhóm. Cô nhìn con bò một lúc, rồi thở dài: ``Nếu chúng ta có chỗ để nhốt nó thì cũng không hẳn là ý tồi.''

``Đến cả bà nữa ư!?'' Pekora và Rushia đồng thanh dầy ngỡ ngàng, gần như ngất xỉu đến nơi.

Đang lúc hỗn loạn, con bò bỗng kêu ``ụm bò\~{}'' một tiếng. Khách hàng xung quanh lấm lét lùi xa, ai cũng hoang mang không hiểu tại sao lại có nguyên con bò đứng chình ình giữa siêu thị. Một bà lão còn nhỏ giọng hỏi nhân viên: ``\vie{Cháu ơi\dots\ ở đây có bán cả bò sống thật à?}''

Người nhân viên đó cũng không biết phải trả lời ra sao. Đến cả bản thân họ cũng không hiểu nổi tại sao lại có một con bò ở bên trong một nơi có không gian hiện đại tới vậy.

Dù rất muốn, nhưng cuối cùng Noel cũng đành từ bỏ ý định mua nguyên con bò. Tuy nhiên, điều đó không có nghĩa là cô nàng sẽ nương tay với các loại thịt.

``\textbf{Tổng tiến công nào mọi người!!!}''

Thế là, cô nàng và đồng bọn quét sạch hai trên ba gian hàng thịt bò, để lại quầy thịt trống trơn với vài khách hàng tội nghiệp đứng đơ người, tay run run cầm giỏ hàng trống rỗng.

``\vie{\dots Tôi chỉ định mua một ít thịt bò cho bữa tối thôi mà\dots}'' một người đàn ông trung niên đứng hình trong tuyệt vọng.

``\vie{H-Họ là ai vậy!?}'' một người phụ nữ khác lẩm bẩm,  nhìn cảnh tượng trước mắt như vừa chứng kiến một cuộc xâm lược của quân Tống.

Flare xoa cằm, nhìn chiếc xe đẩy giờ đã ú ụ thịt, băn khoăn: ``Này, có phải chúng ta hơi quá tay không?''

Noel tuyên bố đầy chắc nịch: `` Không có gì gọi là `quá tay' khi nói về thịt bò!!!''

``Thế còn những người khác thì sao, peko!?'' Pekora chỉ tay về phía các khách hàng đang bàng hoàng.

``Họ có thể ăn thịt lợn hoặc thịt gà mà!'' nàng Hiệp sĩ cười đầy hồn nhiên, như thể là phát súng bắn vào những người đang ngạc nhiên trước một quầy trống trơn kia.

Rushia thì gục xuống, cảm thấy không còn sức lực tranh luận với cô nàng nghiện bò nữa: ``Tôi vái lạy với mức độ cuồng bò của bà rồi đó\dots''

``Vậy nên người ta bảo là `Bò sữa-chan' cũng không sai, peko.''

``Được rồi, được rồi, đi tính tiền thôi\dots\ Trước khi bị bảo vệ đuổi khỏi siêu thị!'' nàng Dị giới thở dài.

Họ cùng nhau kéo Noel đi về phía quầy tính tiền, bỏ lại đám đông khách hàng vẫn còn đang bàng hoàng trước vụ ``cướp sạch'' thịt bò vừa rồi.

\vspace{-40pt}
\subsubsection*{\phantomsection\label{P1.1d}}
\addcontentsline{toc}{subsubsection}{Đội \#{\digi A32\ 192} - Choco, Subaru, Fubuki và Sora}

\begin{center}
    {\color{T4} Nhóm 4 - Choco, Subaru, Fubuki và Sora\\Rau củ}
\end{center}

Cùng lúc này, Choco, Subaru, Fubuki và Sora đang chọn rau củ trong siêu thị cho bữa tối của tiệc sinh nhật lần này.

Họ đứng trước quầy rau củ, nơi bày biện đủ loại rau xanh, củ quả tươi ngon dưới ánh đèn sáng rực. Nhìn những bó rau củ tươi mới và sạch sẽ, những người ăn chay cũng muốn nhỏ nước miếng.

Cô nàng Y tá - không phát cuồng như họ - cũng  là một người rất hứng thú với màu xanh mát của những món thực phẩm giàu chất xơ này.

Cô lướt tay trên những bó rau xanh mướt, ngầm thẩm tính toán như một bà nội trợ chính hiệu: ``Hmm\dots\ Rau nào hợp với ăn lẩu và nướng đây ta\dots''

Bên cạnh, Fubuki cười một cách ranh mãnh: ``Chị cứ chọn rau xà lách làm nướng, còn lẩu thì\dots\ cứ lấy mấy loại rau mà ta hay ăn lẩu ở văn phòng ấy\~{}''

Subaru gật đầu đồng tình, nhưng rồi cũng đưa ra ý tưởng: ``Cũng chưa chắc đâu. Phải thêm cả rau cải và rau muống nữa chứ!''

Choco nghiêng đầu: ``Rau cải thì chị nghe qua rồi, nhưng\dots\ rau muống là gì?''

Sora nhẹ nhàng cười, giải thích: ``Đó là một loại rau đặc sản ở vùng này. Ở chỗ ta bán rất đắt, gấp khoảng mười đến hai mươi lần ở đây cơ.''

Cả ba người còn lại trợn tròn mắt. Subaru há hốc mồm, Fubuki bật cười ``hề hề,'' còn Choco thì giật mình đến mức suýt làm rơi bó rau cầm trên tay.

``S-Sao đắt thế!?'' nàng Vịt trố mắt ngạc nhiên vì sốc nặng; cô không nghĩ rau muống lại có thể cắt cổ như thế.

Nàng Thần tượng cười mỉm: ``Vì chỗ ta không trồng được mà.''

``Em tưởng là nó dễ trồng lắm mà!?''

``Nhưng với khí hậu của nước ta thì chắc nó không trụ nổi đâu\dots'' nàng Cáo trắng đáp lại. ``Xứ mình là xứ cực lạnh mà.''

Cô nàng chắp cằm suy nghĩ, rồi bất chợt nảy ra một ý tưởng: ``Hay là\dots\ tụi mình nhập rau muống về bán với giá cao gấp năm lần? Như thế có khi sẽ thành thương nhân rau muống thành công nhất thế giới đó!''

Choco cũng tỏ vẻ thích thú với ý tưởng điên rồ này: ``Ý hay đó! Cô trò mình chắc sẽ cùng nhau lập công ty chuyên bán rau muống nhập khẩu!''

Thấy thế, nàng Vịt rít lên, muốn dẹp ý tưởng điên rồ đó đi: ``\textbf{Hai người bị làm sao vậy!?}''

Nàng Thần tượng thì chỉ nở một nụ cười hiền, vô tư đáp lại: ``Nếu thế thì để chị làm khách hàng đầu tiên ủng hộ hai đứa nhé\~{}''

``Không phải cả Sora-senpai chứ!!'' Subaru xanh mặt, ngỡ ngàng với lời tuyên bố của cô ấy.

``Chị đùa ấy mà\~{}'' Sora khẽ lè lưỡi, tỏ vẻ ngây thơ.

Nàng Tomboy thấy thế chỉ có thể thở dài bất lực, còn Fubuki và Choco thì hí hửng bàn nhau về kế hoạch kinh doanh rau muống.

\ct

Nhưng rồi họ cũng quay lại chủ đề chính: chọn rau củ cho bữa tiệc sinh nhật tối nay.

``Chúng ta còn phải mua cà chua, khoai, với hành tỏi chứ nhỉ?'' Choco liệt kê. ``Xong rồi còn\dots\ gì nữa ta?''

``Em nghĩ thế là cũng đủ rồi. Cùng lắm chắc mua thêm mấy củ cà rốt cho Pekora-chan nữa thôi,'' Sora đáp lại, rồi sau đó gật đầu khen ngợi: ``Mà, Sensei đúng là tinh ghê!''

``Em nghĩ chúng ta thực sự nên mua một xe cà rốt cho Pekora-chan á\~{}'' Fubuki nhấn mạnh. Trong mừng tượng, cô nghĩ tới viễn cảnh nàng Thỏ đang bơi trong một biển cà rốt, trong khi đang cười một cách đầy điên loạn.

Subaru phủi tay, đập vỡ bong bóng tưởng tượng đó: ``Tôi không nghĩ là em ấy cuồng cà rốt đến vậy đâu. Đừng có mua thừa mứa.''

``Vậy à\dots''

Cả bốn người họ tiếp tục đi dạo một vòng quanh khu rau củ, vừa bàn nhau vừa cân nhắc xem họ nên chọn như thế nào.

Ngoại trừ Fubuki. Cô nàng Cáo trắng này luôn có trò để nghịch ngợm.

Và, lần này cũng không ngoại lệ.

Trong khi ba người họ đang bàn bạc rất hăng say, Fubuki bất ngờ reo lên: ``Subaru, nhìn kìa! Trông nó giống mặt em y hệt luôn!''

``Ý chị là sao chứ?'' Subaru nghiêng đầu khó hiểu.

Nàng Cáo trắng cầm túi cà chua mà cô thấy đó, giơ nó lên cạnh mặt của nàng Vịt vẫn đang ngẩn ngơ kia. Thế rồi, cô cười phá lên: ``Đó! Nhìn hai cái má đỏ đỏ này, y hệt nhau!''

Sora và Choco cũng nhìn quả cà chua và mặt của nàng Vịt, rồi cũng cười mỉm.

\begin{dialogue}
    \speak{Choco} Quả là giống với Subaru-sama nha\~{}
    \speak{Sora} Giống y như đúc luôn á!
\end{dialogue}

Subaru nhăn mặt, gầm gừ: ``Đừng có lấy em ra làm trò đùa nữa!''

Cả nhóm bật cười với phản ứng của cô, nhưng cuối cùng cũng gom đủ rau củ cần thiết.

Riêng Fubuki thì vẫn đang cầm túi cà chua trên tay, lâu lâu lại liếc Subaru với ánh mắt đầy ẩn ý, khiến cho nàng Tomboy không hề thấy thoải mái gì.

\ct

Tại khu vực bày bán trái cây, Choco, Subaru, Fubuki và Sora tiếp tục chọn lựa nguyên liệu cho bữa tiệc sinh nhật.

Trước mặt họ là những loại hoa quả tươi ngon, từ táo, cam, chuối, đến dâu tây, nho, và nhiều loại trái cây khác. Trông loại quả nào cũng tươi và lịm nước cả.

Nàng Y tá chắp cằm suy nghĩ với những loại quả này: ``Hmm\dots\ Ta cần chọn loại trái cây nào để trang trí bánh nhỉ?''

Nàng Vịt cũng toát mồ hôi hột: ``Ca này có vẻ khó đấy. Quả nào trông cũng hợp để trang trí bánh cả\dots''

Nàng Thần tượng nhìn quanh gian hàng một lúc lâu, rồi đưa ra nhận xét: ``Dâu tây là lựa chọn không tệ đâu. Mấy chiếc bánh kem có dâu luôn trông ngon mắt mà.''

Choco gật đầu, cảm thán: ``Thế này thì chuẩn rồi! Vậy ta sẽ lấy chúng nha\dots\ Cùng với cả thêm việt quất, kiwi và đào nữa!''

Nói rồi, hai người họ lấy các loại hoa quả mà nàng Y tá và nàng Thần tượng chọn.

Trong khi đó, Subaru cầm một quả cam lên, bóp nhẹ. ``Còn cái này thì sao? Cam cũng hợp mà.''

Fubuki thấy vậy, thò đầu vào góp ý: ``Cam thì cũng được thôi. Nhưng chị nghĩ là tụi mình cũng nên mua trái cây để ăn sau bữa tiệc nữa! Kiểu như\dots\ dưa hấu chẳng hạn?''

Mắt của nàng Cáo trắng sáng rực khi nhìn thấy một gian hàng chất đầy dưa hấu xanh mướt.

Tất cả mọi người đều cùng đồng tình với ý kiến này. Choco mỉm cười nhẹ: ``Nhắc đến mùa hè thì phải nhắc đến dưa hấu nhỉ?''

``Nhưng mà ta nên lấy bao nhiêu quả đây?'' Sora thắc mắc. ``Chúng ta có tới hơn ba mươi người ăn á\dots''

Fubuki hớn hở cầm một quả dưa hấu to, đáp: ``Em nghĩ là ta cứ lấy thật nhiều!''

Subaru nghe vậy, liền phản biện ngay: ``Không được đâu, Fubuki-senpai! Mua nhiều quá thì ai vác cho hết đây!?'

``Nhưng mà nhưng mà\dots\ nếu chúng ta dùng chúng để chơi trò đập dưa thì sao?''

Tất cả mọi người đứng hình. Nhưng người sốc nhất vẫn là nàng Vịt. Cô không ngờ tới ý tưởng này, một trò chơi mà chỉ có thể thực hiện ở trên biển.

Sunaru há hốc mồm: ``K-Khoan đã, cái gì cơ?''

``Bọn mình bịt mắt lại, cầm một cái gậy, rồi thử đập trúng dưa hấu! Ai đập trúng trước sẽ được ăn phần to hơn, kiểu như vậy!'' nàng Cáo trắng hí hửng giải thích. ``Trò này ta chơi trên biển suốt đấy thôi!''

``Nghe có vẻ thú vị nhỉ?'' Sora vỗ tay, nở một nụ cười đầy nhiệt tình.

Choco thì hơi do dự, lắc đầu: ``Chị thấy trò này hơi nguy hiểm đó, Fubuki-sama. Lần trước chị từng bị văng dưa hấu vào người rồi!''

Subaru lập tức phản đối mạnh mẽ: ``Em cũng không đồng ý đâu! Lần trước em còn bị biến thành dưa hấu vì trò này đấy!''

Fubuki cười hề hề, xua tay: ``Thôi mà, lần này chị sẽ đảm bảo an toàn tuyệt đối! Với lại, có đông người như vậy thì càng vui mà!''

Nàng Y tá nghe vậy, cũng dần bị thuyết phục (một cách quá dễ dàng): ``Ừm\dots\ nếu chuẩn bị kỹ thì chắc cũng không sao đâu nhỉ?''

Nàng Vịt dù vậy vẫn kiên quyết lắc đầu, mặt mày cau có: ``Không! Em không chơi đâu! Em mà lại biến thành dưa hấu nữa thì ai cứu em!?''

Fubuki vỗ vai Choco, cười tươi: ``Dưa hấu thật chứ có phải dưa hấu bị nguyền rủa đâu\~{}''

Choco cũng bật cười, cuối cùng tán thành: ``Thôi được, cô sẽ thử lại lần nữa, nhưng Subaru phải đứng xa ra nhé!'' Cô nàng trước đó đã từng gặp phải một tình huống tương tự như vậy\footnote{Liên hệ đến {\flaregothic ÉPISODE 16}, quyển số Một.} nên lần này nàng Y tá sẽ không để sai lầm mắc lại lần nữa.

Thế là, chỉ có mỗi Subaru là người phản đối kịch liệt nhất, nếu không muốn nói là cảnh giác cao độ nhất: ``Với cả ấy, trò này vốn dĩ phải chơi trên biển mà, với lại lỡ đập hỏng hết dưa thì làm sao ăn!?''

``Cô nhớ hồi em mang dưa hấu tới em cũng đòi chơi trò đập dưa đấy thôi?'' Choco-sensei chỉ ra. ``Sao lần này em phản đối ghê vậy?''

``Vì\dots'' tới đây cô nàng không thể phản bác lại được, vì chính cô là người đã nói như thế từ hai năm về trước, để rồi dẫn đến sự kiện biến thành quả dưa sau vì lời nguyền ``không thể thực thị được ước muốn.''

Nàng Cáo trắng tỉnh bơ đáp lại: ``Em cứ khéo lo! Hỏng thì cứ hốt mấy quả khác ăn thôi\~{}''

``Đừng có bỏ mứa như thế, Fubuki-sama.''

Nàng Vịt tới đây chỉ có thể đập tay vào trán, bó tay với những ý tưởng dị hợm của tiền bối cô. Cô tuyệt vọng: ``Chị nghĩ là tiền mọc ở trên cây hả?''

Fubuki chỉ cười hề hề, trong khi Choco và Sora đã chọn xong vài quả dưa chất lượng, không hề đếm xỉa tới những lời giảng đạo lý của Subaru.

Cuối cùng, sau một hồi tranh luận (mà Subaru là người phản đối mạnh mẽ nhất nhưng vẫn không thắng nổi), họ quyết định lấy vài quả dưa hấu cỡ lớn, trong đó phần lớn là để ăn đàng hoàng, hai tới ba quả để thử trò đập dưa, và số ít quả còn lại dự phòng phòng khi ai đó (có lẽ là Fubuki) lỡ tay làm vỡ mất.

Nhìn một chiếc xe đẩy ú ụ củ quả và hoa quả, nàng Vịt mệt mỏi cất giọng hỏi: ``M-Mấy người đã định dừng lại chưa\dots?'' trong khi hai tay thì nhọc nhằn đẩy.

Nàng Y tá cười tươi với vẻ mãn nguyện: ``Thế này mới đủ cho một bữa tiệc chứ!''

Nàng Cáo trắng thì ôm một quả dưa hấu, cười đầy gian ác: ``Nhớ nha! Chúng ta sẽ đập dưa hấu!''

Nàng Thần tượng giơ nắm đấm đầy nhiệt huyết: ``Vậy thì chúng ta đi thanh toán nào\~{}''

Và thế là bốn người họ kéo nhau về phía quầy tính tiền, mang theo túi lớn túi nhỏ, nổi bật nhất vẫn là những quả dưa hấu tròn căng khiến ai đi ngang qua cũng phải ngoái nhìn.

Kỳ quặc thay, đây là nhóm đi mua sắm yên bình trong tất cả các nhóm. Ngoại trừ việc họ mua với số lượng lớn khiến người thu ngân phải toát mồ hôi hột, còn lại thì còn trông như những người đi sắm bình thường khác.

\vspace{-40pt}
\subsubsection*{\phantomsection\label{P1.1e}}
\addcontentsline{toc}{subsubsection}{Đội \#{\digi\ 100698} - Suisei, Marine, Polka, Kanata, Towa, Watame, Luna và Nene}

\begin{center}
    {\color{T5} Nhóm 5 - Suisei, Marine, Polka, Kanata, Towa, Watame, Luna và Nene\\Quà tặng}
\end{center}

Suisei, Marine và Polka lang thang trong khu vực quà tặng, cân nhắc xem nên mua gì cho cậu quản lý. Mặc dù Kuon - hay theo lời Ryuuhou - có nói rằng ``không cần quà cáp gì'', nhưng cả ba vẫn muốn chọn một món quà thật ý nghĩa để thể hiện lòng thành của mình.

Polka cầm lên một bộ bài Tây, giơ ra trước mặt hai người còn lại, hớn hở hỏi: ``Cái này có được không?''

Nhìn thứ mà nàng Cáo sa mạc cầm trên tay, cả nàng Thuyền trưởng và nàng Ca sĩ đều cảm thấy dòng điện chạy dọc sống lưng. Cả hai người họ run lẩy bẩy, sợ sệt khi nhớ lại khoảnh khắc đó.

``E-Em định khơi mào thêm một trận Tiềm Ô Quân nữa ư\dots!?'' Marine hoảng hốt, mồ hôi túa lúa như suối.

``Em chưa biết là nó kinh khủng đến cỡ nào đâu\dots'' mặt của Suisei tối sầm lại, toát ra một thứ sát khí khiến cho khách hàng quanh cô phải tránh sang một bên.

Polka nhìn họ với ánh mắt ngơ ngác đầy khó hiểu: ``Ể? Hai chị bị sao vậy? Em định mua nó về để làm ảo thuật mà.''

Cả hai cô gái đồng loạt thở phào nhẹ nhõm. Nhưng Marine vẫn thắc mắc: ``Cái đó là mua cho em, nhưng mà ta đang mua cho Kuon-senpai mà?''

Nàng Chủ xiếc cười tươi đáp lại: ``Thì đúng, nhưng em muốn diễn trò ảo thuật cho mọi người ở nhà anh ấy xem!''

Nghe cũng có lý, nhưng có vẻ vẫn chưa phải thứ thực sự dành cho cậu Tổng của họ.

Lúc này, Marine bỗng giơ một cái chai lên đầy tự đắc: ``Sao ta không tặng một cái chai nhỉ?''

Suisei và Polka nghiêng đầu nhìn cô: ``Ủa? Làm gì?''

``Để anh ấy chơi trò lắp mô hình thuyền trong chai, giống của chứ sao!''

Nàng Cáo xiếc nhìn cô với vẻ nghi hoặc. Cô thở dài, khoanh tay đáp lại với chút châm chọc: ``Có phải ai cũng có sở thích cổ lỗ sĩ như chị đâu\dots''

Tới đây, Senchou bùng nổ, cảm tưởng như bị xúc phạm: ``Này, này, ý em cổ lỗ sĩ là sao chứ?''

Suisei mặc kệ màn tranh cãi của hai người kia, cô bỗng nghĩ ra một ý tưởng khác: ``Sao ta không tặng cho anh ấy một thứ mà anh ấy thích nhất nhỉ? Kiểu, cứ làm theo gợi ý của Ryuuhou-chan ấy?''

Marine và Polka dừng tranh luận, quay sang nhìn cô.

Nàng Thuyền trưởng: ``Cũng có thể. Cơ mà em ấy nói gì thế nhỉ\dots''

Nàng Cáo xiếc tiếp tục trêu chọc: ``Vừa mới nói hồi sáng xong mà lại quên nhanh thế, bà già?''

``Chị đánh chết em bây giờ,'' gân máu của Senchou hiện ra, đôi mắt cô nheo lại đầy tức giận. Đáp lại, Polka chỉ cười một cách khoái chí.

Nàng Ca sĩ mỉm cười, nói với một chút hào hứng: ``Em ấy nói là Kuon-senpai muốn có một chiếc S*itch!''

``À\dots\ ra thế,'' Marine thở dài. ``Đúng là Ryuuhou-chan nói như vậy thật.''

``Cái đó thì dễ thôi!'' Polka đáp lại đầy nhiệt tình. ``Ta tới gian điện tử mà lựa thôi!''

Cả ba phấn khích tiến đến khu vực bán đồ điện tử, bắt đầu săn hàng.

\ct

Khu vực của các máy chơi điện tử là một trong những nơi được bày bán những chiếc máy chơi hiện đại và tân tiến nhất mà Nhật Bản có (vì đây là AEON mà, nhớ chứ?). Về cơ bản, ở xứ sở mặt trời mọc có gì, thì ở đây cũng có. Điều đó khiến cho các cô gái \hll\ dễ dàng chọn lựa hơn.

Trước mặt ba cô gái, một dãy S*itch với nhiều phiên bản và màu sắc khác nhau được sắp xếp gọn gàng, đi cùng với đó là các kệ trò chơi đi kèm.

Marine xoa cằm, tỏ vẻ lưỡng lự: ``Nên chọn bản thường hay OLED đây?''

``Tất nhiên là bản OLED rồi! Màn hình to đẹp hơn, pin trâu hơn!'' Suisei nghiêm túc nói, đáp một cách gọn ghẽ.

``Nhưng mà màu nào đây?'' Polka tiếp tục hỏi.

Họ nhìn lướt qua các mẫu: từ các màu cơ bản như xanh - đỏ, xám, trắng, cho tới những bản theo chủ đề như Z*lda và Ma**o.

Trước những chiếc máy chơi game nhà Nin, Cáo Sa mạc đưa ra giải pháp được coi là an toàn nhất: ``Em nghi ta cứ chọn hai màu vàng-đỏ đi? Vừa có màu mà anh ấy yêu thích, vừa hợp với tên của anh ấy nữa!\footnote{Giải thích: chữ ``kawa'' có hai cách đọc: 鈹 và 川, trong đó gốc  của chữ thứ nhất có nghĩa là ``màu vàng''.}''

Cả nàng Ca sĩ và nàng Thuyền trưởng bật cười với lựa chọn này. ``Thế cũng được đấy chứ!''

Họ quyết định sẽ theo ý tưởng của Polka. Đồng thời, để giữ bí mật, họ còn yêu cầu nhân viên gói lại một cách gọn gàng.

\ct

Sau khi chọn xong máy, họ tiếp tục bàn bạc về game.

Marine xoa cằm, mắt quét qua dãy game xếp ngay ngắn trên kệ. Cô nàng hửng: ``Ta có nên mua kèm game luôn không?''

Suisei gật đầu đồng tình: ``Đương nhiên rồi! Nhưng mà\dots\ nên mua game gì ta?''

Cả ba cô gái bắt đầu liệt kê ra những trò chơi mà họ có thể nghĩ được.

``Ma**o Kart này!'' Polka nhí nhảnh với đôi mắt sáng rực.

``Sm*sh B**s nữa!'' Marine theo sau hiến kế, khúc khích cười.

Suisei thì thậm chí còn tỏ vẻ nguy hiểm hơn: ``Hoặc Taiko no Tatsujin\dots\ để ép anh ấy chơi trống!?''

Cả ba nhìn nhau với ánh mắt đầy toan tính, nhưng chưa dừng lại ở đó. Nàng Ca sĩ tiếp tục lướt mắt qua kệ game, rồi đột nhiên cầm lên một hộp trò chơi quen thuộc.

``Còn Tetris 99 thì sao?'' cô nàng tưoi cười, tay cầm một trong những trò quen thuộc.

Vừa dứt lời, cả Marine lẫn Polka đồng loạt biến sắc.

``K-khoan đã, Suisei-senpai\dots\'' Polka tái mét. ``T-Tetris á!?''

Marine cũng hoảng hốt không kém: ``C-Chị định bắt Kuon-senpai chơi với chị trò đó á!?

Suisei nhìn họ, vẻ ngây thơ vô (số) tội. Cô chớp mắt đầy ngạc nhiên: ``Thì có sao đâu? Trò này cũng hay mà.''

Nàng Cáo sa mạc rùng mình: ``Không không không! Cái này mà vào tay chị thì chả khác nào dụ một con cừu non vào hang hổ cả!''

Senchou cũng gật đầu lia lịa: ``Đúng đó! Để Kuon-senpai đối đầu với chị trong Tetris chẳng khác nào bắt anh ấy trèo lên đỉnh Everest bằng một chân!''

Họ tiếp tục đưa ra bằng chứng lập luận, cố gắng thuyết phục nàng ``Sao chổi''  từ bỏ ý định mà họ cho là ngớ ngẩn này. Nhưng đáp lại, Suisei chỉ bật cười trước phản ứng thái quá của họ.

``Hai người làm quá rồi. Chị với Kuon-senpai từng đấu giao hữu trước đây rồi\footnote{Xem lại phần {\abv Truyện phụ}, {\flaregothic ÉPISODE 44}, quyển số Bốn.}. Chị biết khả năng của anh ấy đến đâu mà.''

Polka và Marine nhìn nhau, vẫn có chút bất an.

Nàng Cáo xiếc ngập ngừng; ``Nhưng mà\dots\ chắc chắn là chị không có đè bẹp Kuon-senpai đấy chứ?''

``Không phải kiểu vừa vào game là ảnh bị chị đánh úp bằng một chuỗi combo hủy diệt rồi tắt điện luôn!?'' Marine trợn mắt, mặt trắng bệch ra.

Nàng Ca sĩ chỉ nhún vai đáp lại: ``À thì\dots\ có một chút.''

``Một chút\dots\!?'' cả hai người thất thần, nghĩ về cảnh tượng không mấy tốt đẹp tới cậu Tổng.

Nhưng đến cuối cùng, Suisei chỉ nhìn hai cô gái hoảng hồn kia một cách thản nhiên, rồi đặt hộp game vào giỏ.

``Hy vọng là anh ấy bình an trước chị ấy\dots\'' họ lẩm bẩm, mắt ngước lên trần nhà như đang cầu xin Thần linh có thể nghe được tiếng lòng của mình.

\ct

Sau khi ổn định tinh thần, cả ba tiếp tục lùng sục thêm đĩa trò chơi. Họ bàn bạc rôm rả đến mức khiến khách hàng xung quanh phải ngoái lại nhìn. Trong mắt họ, ba cô gái trông chẳng khác những tay chơi game lão luyện là bao.

Marine phấn khích cầm một đĩa khác: ``Này, còn Splatoon thì sao? Chơi theo nhóm với mọi người vui lắm!''

Polka đập tay, tỏ vẻ đồng tình: ``Được đó! Nhưng Kuon-senpai có thích bắn súng sơn không?''

Suisei khoanh tay, trầm ngâm ngẫm nghĩ: ``Hừm\dots\ anh ấy không ghét đâu, nhưng cũng không phải kiểu cuồng nhiệt gì lắm.

``Hình như anh ấy cũng không phải kiểu người thích RPG lắm nhỉ?'' nàng Cáo sa mạc bổ sung.

``Anh ấy lại thích những trò giải đố hơn\dots\'' Senchou cũng gật đầu.

Cả ba cô nàng đứng đó suy nghĩ một hồi lâu, quyết định xem liệu họ có lấy chiếc đĩa này không.

``Cơ mà nếu là để chơi cùng mọi người thì chắc anh ấy cũng sẽ thử thôi,'' nàng Ca sĩ nhún vai.

Thế là, cả ba người quyết định đặt đĩa game đó vào giỏ hàng.

Ở gần đó, một vị khách lén lút rời khỏi kệ game, lắc đầu ngao ngán khi nghe ba cô nàng này bàn luận một cách nghiêm túc đến mức đáng sợ về sở thích chơi điện tử của cậu Tổng quản lý.

Dù thế, cả ba cô gái vẫn mặc kệ những ánh mắt dò xét từ mọi ngưòi, tiếp tục bàn tán một cách nhiệt tình.

``Này này! Mua thêm Ring Fit để ép anh ấy tập thể dục không?'' nàng Cáo sa mạc hớn hở.

Senchou gật đầu cái rụp: ``Ý hay, ý hay! Đỡ bị còng lưng do làm việc văn phòng nhiều quá!''

Suisei thì mỉm cười gian xảo: ``Và cũ g là một cái cớ để bắt Senpai mặc đồ thể thao\~{}''

Polka và Marine bật cười khúc khích, rồi nhanh tay vớt luôn hộp game Ring Fit vào giỏ.

Đến khi họ nhìn lại, giỏ hàng đã đầy ắp. Đếm sơ sơ thì họ đã lấy gần hết bộ sưu tập trò chơi mà siêu thị này có.

Marine há hốc miệng trước số lượng trò chơi mà đã lấy: ``Chúng ta có\dots\ hơi lố không nhỉ?''

Polka chỉ bĩu môi đáp, chẳng hề để tâm: ``Sinh nhật mà! Chẳng lẽ keo kiệt với anh ấy được sao?''

Suisei cười một cách ranh mãnh, tỏ vẻ hài lòng: ``Chị nghĩ là không đâu.''

Sau khi chọn xong, họ mang cả đống game lẫn chiếc máy chơi điện tử đến quầy thu ngân. Nhìn số lượng đĩa trò chơi mà họ đã lấy, nhân viên đứng quầy cũng ngạc nhiên, tới mức không biết nên vui vì ``thần tài tới'', hay nên sốc vì họ mua số lượng lớn nữa.

Nàng Ca sĩ mỉm cười yêu cầu; ``Nhờ bạn gói thành quà giúp bọn mình nhé!''

Nhân viên cũng chỉ cười gượng đáp lại: ``Dạ\dots\ để mình gói cho các bạn ngay.''

Người nhân viên nhìn thấy số lượng game thì có hơi sững sờ, nhưng vẫn nhanh chóng làm việc của mình. Trong khi đó, ba cô nàng đứng bên cạnh, khoanh tay đầy tự mãn vì đã hoàn thành nhiệm vụ một cách xuất sắc.

Polka hí hửng cười tít mắt: ``Chắc chắn đây là món quà khủng nhất trong tất cả luôn rồi!''

Marine gật gù, cười thỏa mãn: ``Ừm! Anh ấy mà không vui thì ta bóp cổ cho vui luôn!''

``Để xem Kuon-senpai phản ứng ra sao  khi mở ra nào\~{}'' Suisei nở một nụ cưòi thắm thiết, nhưng ẩn sâu trong đó có một chút nguy hiểm.

Và thế là, với một túi quà sinh nhật siêu nặng đô, họ hớn hở rời khỏi quầy, thỏa mãn với chiến lợi phẩm của mình.

\begin{center}
    \uline{Quầy đồ thủ công.}
\end{center}

Ở một diễn biến khác\dots

Bộ tứ của Thế hệ thứ tư, cùng với Nene, đang lang thang giữa các quầy hàng đồ thủ công, tìm kiếm những món quà sinh nhật nhỏ gọn hơn dành cho Kuon. Nhưng không chỉ vậy, họ còn tranh thủ tìm quà cho chính mình và cả các thành viên khác làm kỷ niệm.

Trước những phụ kiện nhỏ nhắn này, Towa nhíu mày đầy suy tư: ``Không biết anh ấy thích gì nhỉ\dots''

Kanata cầm một vòng tay bằng chun lên, nhún vai: ``Tôi nghĩ ta cứ tìm những món nào phù hợp với tính cách của anh ấy thôi.''

Luna, vốn đang nghịch một chiếc dây đeo tay bằng da, bỗng reo lên vẻ thích thú: ``Zậy dây đeo và vòm cổ tay có đực khum na\~{}''

Nàng Thiên sứ cầm thử một cái lên, săm soi một lúc rồi gật gù ưng: ``Cũng được. Nhưng còn tùy xem liệu Kuon-senpai có hợp không nữa.''

``Luna nghĩ là anh í xẽ thít thui-nora\~{}'' nàng Công chúa xứ Kẹo ngân nga, tay lắc một chiếc vòng bạc khác.

Towa nhíu mày, tỏ vẻ thận trọng: ``Tôi cũng không biết nữa\dots\ Hơi căng đấy\dots''

Trong khi ba người còn đang lăn tăn về sở thích của cậu Tổng, Watame và Nene đã rẽ sang quầy phụ kiện gần đó, tay cầm đủ loại vòng tay khác nhau.

Nàng Cừu giơ lên một chiếc vòng, hỏi cô nàng Hải cẩu: ``Em thấy cái vòng tay màu cam này thế nào?''

Nàng Hải cẩu nheo mắt bình phẩm: ``Có thêm hình củ cà rốt nữa thì tốt nhỉ-aru\dots''

``Ừm, vậy chắc là cái vòng này?'' Watame cầm một chiếc khác giơ lên, lần này đính kèm một củ cà rốt.

``Sáng hơn một chút được không nhỉ?''

Và thế là, hai người cứ thế tìm kiếm, thử hết cái này đến cái khác, bàn bạc rôm rả như đang lựa đồ cho chính mình vậy.

Mười phút trôi qua. Sau khi hai người họ bàn tới bàn lui, cuối cùng Nene đã có chọn được chiếc vòng tay ưng nhất cho mình.

``Đây rồi! Chiếc này là chuẩn nhất!'' cô nàng nở một nụ cười mãn nguyện.

``Vậy là tìm được rồi nha\~{}''

Cả hai vỗ tay ăn mừng nho nhỏ, ôm chặt lấy ``chiến lợi phẩm'' trong tay như thể vừa đoạt được kho báu.

\ct

Trong khi Watame và Nene còn đang hí hửng với vòng tay của mình, ba người còn lại cũng đã chọn được một vài món, nhưng thay vì nghiêm túc chọn cho Kuon, họ lại đang lăn ra\dots\ châm chọc nhau.

Luna cầm một chiếc vòng cổ đính đá nhỏ lấp lánh, mắt sáng rực: ``Towapi, cái nì hợp zới bà lứm đó-nora\~{}!''

Towa nhăn mặt, đổ mồ hôi hột: ``Trông nó lòe loẹt quá trời! Cái này phải hợp với ai thích lộng lẫy như Marine-senpai mới đúng chứ?''

``Hoặc là Polka cũng được\dots'' Kanata khúc khích. ``Em ấy mê mấy món như vậy đó.''

Nàng Công chúa xứ Kẹo đưa một cái dây chuyền bạc đơn giản hơn cho hai người: ``Thế còn cái này? Hợp với bà mừ-nora\~{}?''

Nàng Ác ma cầm thử lên, soi xét một hồi: ``Ừm\dots\ cũng được\dots'', nhưng rồi, chợt nhíu mày khi sực nhớ ra điều quan trọng: ``Đợi đã, này đâu phải cho tôi!''

Nàng Thiên sứ giả vờ thở dài: ``Vậy rốt cuộc, chúng ta đang mua cho Kuon-senpi hay mua cho chính mình thế?''

``Chắc là cả hai?'' cả hai người họ đồng thanh, với một giọng lưỡng lự, một giọng tỉnh bơ.

Luna bưng miệng cười: ``Tui nghĩ tui xẽ mua một cái cho Choco-sensei á-nanora\~{}''

``Cái nào?''

Nàng Kẹo chỉ vào một sợi dây có mặt hình trái tim. ``Tui nghĩ cô ấy xẽ thích lắm á!''

Towa gật đầu đồng ý. Cô giơ lên một cái vòng cổ có mặt hình con bò sữa lên, trề môi: ``Cái này thì chắc chắn hợp với Noel-senpai này.''

Cả ba người họ cười phá lên, trong đầu tưởng tượng tới viễn cảnh nàng Hiệp sĩ đang cải trang thành một con bò sữa.

Kanata cũng không chịu thua: ``Nhưng nếu vậy thì cái này phải dành cho Fubuki mới đúng!'' rồi giơ lên một dây chuyền có mặt hình măng non.

``Còn cái này chắc chắn là dành cho Miko-chi rồi\dots'' Towa giơ lên một cái vòng cổ hình cây gậy trừ tà.

Ba người cứ thế chuyền qua chuyền lại, bình luận đủ kiểu về từng chiếc vòng cổ và ai sẽ hợp với cái nào, đôi khi còn cười đùa đến mức suýt làm rơi đồ.

Nhưng cuối cùng, sau một hồi bàn bạc và chọc nhau không chán, họ cũng chọn được những món vừa ý cho các thành viên \hll\ và cho chính Kuon: một cái vòng đeo tay đơn giản màu đỏ.

``Được rồi, giờ thì tới quầy tính tiền thôi!'' nàng Thiên sứ - giờ làm Trưởng nhóm của Thế hệ thứ Tư giơ nắm đấm lên trời.

``\textbf{Oooh\~{}!!}''

Dứt câu, họ tụ hội lại với nhóm tặng quà đã đến trước đó là Suisei, Marine và Polka tại quầy thanh toán.

\vspace{-40pt}
\subsubsection*{\phantomsection\label{P1.1f}}
\addcontentsline{toc}{subsubsection}{Đội \#{\digi A65402} - Mio, Miko, Matsuri, Aki, Mel và Botan}

\begin{center}
    {\color{T6} Nhóm 6 - Mio, Miko, Matsuri, Aki, Mel và Botan\\Đồ trang trí}
\end{center}

Mio và Miko đang đi dạo quanh khu vực bán đồ trang trí, cố gắng tìm ra cách tốt nhất để khiến phòng tiệc trở nên lộng lẫy nhất có thể.

Mỗi tội, phong cách trang trí của hai cô nàng lại có phần trái ngược, thậm chí là xung khắc nhau.

Nàng Vu nữ hào hứng chỉ vào một gói đồ được bọc rất cẩn thận, trên đó có ghi hai từ: ``Đá khô''. ``Đá khô có được không ne?''

Nàng Sói đen nhìn chằm chằm vào nó, rồi quay sang hỏi: ``Chị định dùng nó làm gì vậy?''

Miko vung tay đầy nhiệt huyết, hào hứng: ``Miko xem trên mạng thấy người ta dùng nó để tỏa khói ấy, nhìn đẹp cực đó ne! Kiểu như bước vào phòng tiệc mà có làn khói nhẹ trôi trên sàn ấy, cứ như bước vào thế giới phép thuật ne\~{}!''

Mio nghe vậy, liền đổ mồ hôi hột: ``Nghe nguy hiểm quá đó! Nếu mà lỡ ai đụng vào rồi bị bỏng lạnh thì sao!?''

``Thế thì ta chỉ cần cẩn thận là được ne\~{}''

Dù vậy, nàng Sói đen vẫn lắc đầu kiên quyết: ``Không được! Ta trang trí phòng cho tiệc chứ đâu phải định lập phòng thí nghiệm hóa học đâu!''

Đến đây, nàng Vu nữ xụ mặt, phồng má: ``Mio-chan đúng là mất vui\dots\ Nếu thế thì ta trang trí kiểu gì đây ne\dots?''

Mio khoanh tay lại, suy nghĩ một chút để tính toán. Cô đưa ra một gợi ý thay thế phù hợp hơn: ``Sao ta không dùng ruy băng, bong bóng với chữ chúc mừng sinh nhật ấy?''

``Thế thì chán lắm ne!'' Miko bất mãn.

``Nhưng ít nhất nó còn an toàn hơn\dots'' Mio thở dài, cố gắng dập tắt những ý tưởng điên rồ từ tiền bối cô.

Tất nhiên, nàng Vu nữ vẫn chưa chịu từ bỏ. Cô tiếp tục nghĩ ra một ý tưởng khác: ``Hay là\dots\ pháo giấy? Hoặc là lắp nguyên cái màn hình LED để chạy hình ảnh của mọi người? Hay thậm chí thuê luôn một dàn đèn sân khấu, ne!?''

Những ý tưởng mà người thường không nghĩ được ấy làm cho nàng Sói đen sốc toàn tập. ``Chúng ta đang tổ chức sinh nhật trong nhà người ta, không phải mở concert đâu!''

``Thế ít nhất cũng phải có thứ gì đó hoành tráng chứ, ne!''

Mio xoa thái dương, bắt đầu nhượng bộ: ``Thôi được rồi, ta có thể thêm vài cái đèn nhấp nháy và treo một số đồ trang trí có phát sáng. Nhưng không có chuyện dùng đá khô hay pháo giấy đâu!''

``Ể\~{}?'' Miko chớp mắt. ``Nhưng mà pháo giấy vui mà!''

Tới đây, nàng Sói đen bắt đầu đánh vào tâm lý lười biếng của đối phương: ``Bộ chị muốn dọn dẹp nhà anh ấy sau bữa tiệc ấy hả?''

Và, đúng như cô dự đoán. Nàng Vu nữ lập tức lắc đầu nguầy nguậy như cái chong chóng: ``Không! Không đâu, ne!''

``Vậy thì bỏ mấy ý tưởng này đi nhỉ?''

Sau một hồi giằng co giữa sự điên rồ của Miko và sự thực tế của Mio, cuối cùng họ cũng thống nhất với nhau một số ý tưởng hợp lý hơn. Miko vẫn chưa hoàn toàn từ bỏ ý định làm gì đó đặc biệt, nhưng trước mắt, họ quyết định chọn những món đồ trang trí phù hợp trước.

\ct

Miko đang bước qua từng gian hàng, mắt đảo quanh như thể đang tìm kiếm điều gì đó thú vị hơn những đồ trang trí thông thường. Khi đi ngang qua một dãy kệ chất đầy thùng sơn đủ màu sắc, mắt cô nàng bỗng sáng rực lên như vừa nảy ra một ý tưởng thiên tài.

``\textbf{Mio-chan\~{}!!} Ta vẽ tranh sơn tường đi ne!!'' cô nàng reo lên như một đứa trẻ.

Mio khựng lại với tiếng reo đầy bất ngờ của tiền bố mình.

``Nhìn xem nè, có đủ màu luôn! Nếu ta vẽ một bức tranh lớn trên tường phòng tiệc thì chắc chắn sẽ hoành tráng hơn nhiều so với mấy cái ruy băng đơn giản à!''

``Nhưng mà\dots\ đó là nhà của người ta mà! Không phải của chị đâu!'' nàng Sói đen nhăn mày, phản bác lại.

``Nhưng mà nhé, Korone-chan và Okayu-chan đã từng sơn tường nhà Kuon-senpai từ Tết năm ngoái\footnote{Xem lại {\sen Partie 2}, quyển T.} rồi còn gì\~{}'' nàng Vu nữ tiếp tục dùng giọng nói ngon ngọt của mình để thuyết phục cô nàng Sói đen kiên định ấy.

``Thế thì liên quan gì?''

``Thì ít nhất, ta biết là anh ấy không phản đối chuyện thay đổi bức tường mà!''

Mio thở dài, giải thích: ``Nhưng sơn tường khác hẳn với việc vẽ hẳn một bức tranh lên đó! Nếu gia đình anh ấy không thích thì sao?''

Dù thế, nàng Vu nữ nở một nụ cười ranh mãnh. Cô tiếp tục xoáy sâu vào cảm xúc của cậu Tổng năm khi nhìn thấy bức tường được sơn năm đó - điều mà cậu cũng không ngờ tới: ``Nhưng mà nè\~{} Mio-chan nhớ không? Lần đó khi hai nhóc ấy sơn lại tường, Kuon-senpai đã nhìn nó một lúc, chỉ khẽ ngạc nhiên rồi thậm chí xoa đầu hai em ấy, chẳng hề phản đối hay trách mắng gì cả ne!''

Tới đây, nàng Sói đen bắt đầu lung lay lập trường của mình: ``Ừ thì\dots\ Đúng là vậy, nhưng mà\dots''

Thấy cô đã bắt đầu cảm thấy khó xử, Miko tiếp tục nhấn mạnh: ``Mà anh ấy cũng thuộc dạng dễ tính mà ne! Nếu bọn mình vẽ đẹp thì chắc chắn sẽ ổn thôi!''

``Nhưng mà\dots\ nếu vẽ xấu thì sao?''

``Không lo! Miko đã luyện vẽ suốt bao năm qua rồi ne!''

``Thế lần cuối chị vẽ là khi nào?'' Mio liếc nhìn cô với đầy sự nghi ngờ.

Nàng Vu nữ gãi đầu gãi tai, cười khì: ``Từ hồi\dots\ còn đi học?''

Nàng Sói đen chỉ có thể thở dài với câu trả lời đầy ngây thơ từ tiền bối mình.

Hiển nhiên, Miko vẫn chưa chịu thua. Cô thậm chí còn kéo tay cô nàng Sói đen, tiếp tục nói những lời ngon ngọt với đầy sự quyết tâm: ``Đi mà Mio-chan! Chị cũng thích làm gì đó đặc biệt đúng không ne\~{}? Chị dám chắc là các thành viên khác cũng muốn như vậy lắm á!''

Nàng Sói đen cũng chỉ có thể bật cười đầy bất lực, đành giơ cờ trắng đầu hàng: ``Thôi, thôi được rồi. Chúng ta cứ hỏi ý kiến mọi người đã.''

``\textbf{Quá đã ne!!}'' nàng Vu nữ giơ hai tay lên trời đầy phấn khích.

Vậy là, ngoài việc chỉ đơn thuần trang trí bằng ruy băng và bong bóng, bộ đôi này quyết định sẽ vẽ hẳn một bức tranh lên tường phòng tiệc. Tất nhiên không chỉ hai người họ, mà còn các thành viên \hll\ khác nữa.

Một quyết định có phần điên rồ, nhưng lần này, ngay cả Mio cũng phải gật đầu đồng ý (dù có chút miễn cưỡng).

\ct

Matsuri, Aki, Mel và Botan đang đi qua từng quầy hàng, tìm kiếm những món đồ trang trí phù hợp cho bữa tiệc sinh nhật. Họ đã chọn được một số thứ cơ bản như bong bóng, dải ruy băng và đèn nhấp nháy, nhưng nàng Cổ động dường như vẫn chưa hoàn toàn hài lòng. Cô tiếp tục đi loanh quanh, mắt lấp lánh như trẻ con vào dịp lễ hội, rồi bất chợt cầm lên một cái nơ to màu đỏ và cắm thử lên đầu mình.

Cô xoay một vòng đầy điệu nghệ, cười tươi roi rói: ``Mọi người thấy mình thế nào\~{}?''

Aki và Mel lập tức vỗ tay hưởng ứng.

Nàng Tóc bay cười dịu dàng, vỗ tay nhiệt tình nhất: ``Dễ thương quá chừng luôn\~{}!''

Nàng Dơi gật đầu đồng tình: ``Trông cứ như tiểu thư của một buổi tiệc quý tộc vậy!''

Nghe hai lời khen của hai người cùng Thế hệ, cô cười tít mắt: ``Đúng không nào\~{} Mình nghĩ có khi sẽ đội nó luôn đến bữa tiệc ấy!''

Trong khi ba người đang ríu rít khen ngợi nhau, Botan vẫn im lặng đứng một chỗ, ánh mắt nhìn chằm chằm vào giá trưng bày phía trước như thể đang suy nghĩ gì đó rất nghiêm túc.

Botan không hề tham gia vào màn tán dương này. Nàng sư tử đứng yên tại chỗ, mặt cứng đờ,

Thấy vậy, Mel nghiêng đầu đầy khó hiểu: ``Botan-chan? Sao thế em?''

Aki thì vẫy tay trước mặt: ``Em ổn chứ?'' nhưng không nhận được câu trả lời từ cô ấy.

Không nói gì, Botan từ từ đưa tay lên, cầm lấy một chiếc bờm có gắn hai khẩu súng nhỏ hai bên - một món đồ trang trí khá\dots\ lạ so với những thứ xung quanh.

Cô nhìn chằm chằm vào nó trong vài giây, rồi như bị hấp dẫn bởi một sức mạnh vô hình, nàng Sư tử lặng lẽ đeo nó lên đầu, rồi - trong một khoảnh khắc hiếm hoi - môi cô khẽ cong lên thành một nụ cười thích thú.

Ngay lập tức, cả ba người còn lại đều tròn mắt nhìn, rồi đồng loạt cười rộ lên.

Matsuri đập tay vào vai cô, bụm miệng cười: ``Ủa trời!? Em thích cái đó lắm đúng không?''

Aki thì tinh tế che miệng cười: ``Không ngờ lại hợp với em luôn đó!''

Mel thì cười một cách tinh nghịch nhưng cũng không giấu nổi vẻ thích thú: ``Giống như `Sư Tử Chiến Binh' ấy nha\~{}!''

Nghe những tiếng cười của đàn chị, Botan giật mình, luống cuống gỡ bờm của mình xuống, chống chế: ``K-Không phải như các chị nghĩ đâu! Không phải! Chỉ là\dots\ nhìn nó thú vị nên em thử thôi!''

Nhưng thay vì để cô lảng tránh, ba bậc tiền bối chỉ càng cười lớn hơn, còn cố tình nhìn nhau gật gù ra vẻ hiểu rõ.

``Vậy sao em cứ tủm tỉm cười khi soi gương thế?'' nàng Cổ động cười một cách ranh mãnh, rồi rướn đến sát mặt của đàn em.

Nàng Sư tử quay mặt đi, lúng túng: ``T-Tại\dots\ Tại vì\dots''

Nhưng chưa kịp nghĩ ra lý do hợp lý, cô đã bị ba người còn lại cười châm chọc, vây quanh khen ngợi tới tấp. Botan cố giữ vẻ mặt bình tĩnh, nhưng đôi tai đã đỏ lựng như thể sắp bốc khói.

Cuối cùng, cô bực mình hất tay: ``Thôi, thôi, lo chọn đồ đi, đừng có trêu em nữa\dots\ Các senpai đúng là quá đáng\dots''

``Vậy là em thừa nhận rồi ha!'' Matsuri cười đầy tinh quái, khiến cho nàng Sư tử chỉ có thể thở dài đầy bất lực.

Dù có hơi ngại ngùng, nhưng cuối cùng Botan cũng không tháo chiếc bờm xuống ngay mà tiếp tục đeo nó trong lúc đi lựa đồ trang trí. Cô không nói ra, nhưng có vẻ như cô cũng khá ưng ý với món đồ này.

Nhưng tất nhiên, Botan quyết định sẽ không để yên chuyện này. Cô nheo mắt, khóe môi khẽ nhếch lên, trong đầu nảy ra một kế hoạch ``trả thù'' nhẹ nhàng.

Cô giả vờ lơ đãng, cất giọng hỏi: ``Mà nè\dots\ chẳng phải bọn chị cũng nên thử vài món đồ hợp với phong cách của mình sao?''

``Hửm? Ý em là sao?'' nàng Cổ động chớp mắt.

Botan không đáp lại, mà chỉ lẳng lặng với tay lấy một chiếc bờm tai mèo màu hồng lấp lánh và dúi thẳng lên đầu cô.

Matsuri bất ngờ với hành động kỳ quặc của cô: ``Ể!? C-Cái gì đây!?''

Nàng Sư tử không để tâm, chỉ ngắm nghía tiền bối một hồi rồi gật gù: ``Ừm\dots\ Hợp thật đấy.''

Lập tức, cả Aki và Mel đều cùng hồ hởi với khuôn mặt bầu bĩnh của Matsuri: vừa có bờm tai mèo, lại có chiếc nơ đỏ cắm trên đầu nưa.

\begin{dialogue}
    \speak{Mel} Matsuri-chan dễ thương thật đó!
    \speak{Aki} Đúng rồi đó! Matsuri-chan đội tai mèo dễ thương lắm!
\end{dialogue}

Nghe những lời có cánh từ hai người, nàng Lễ hội đỏ mặt, cố gỡ chiếc bờm xuống: ``T-Thôi nào, bà đừng có mà-''

Nhưng chưa kịp làm gì, Botan đã nhanh chóng lấy thêm một chiếc bờm khác - lần này là loại có sừng quỷ - và nhẹ nhàng cài lên tóc Mel.

``E-Ê!?''

``Mel-senpai có vẻ hợp với phong cách huyền bí này hơn,'' nàng Sư tử mỉm cười.

Nàng Tóc bay gật đầu hưởng ứng: ``Ừ nhỉ! Nhìn như một nữ hoàng quỷ quyến rũ ấy!''

Nàng Dơi nghe vậy, liền che mặt đầy ngượng ngùng: ``C-Cái gì mà nữ hoàng chứ\~{}! Đừng có mà trêu tôi nữa đi\~{}!''

``Ôi trời,'' nàng Lễ hội cười gian, ``nhưng mà công nhận nhìn hợp phết đấy nha\~{}''

``Đừng có nhìn tôi như vậy mà\dots''

Chưa dừng lại ở đó, Botan tiếp tục nhặt lên một chiếc mũ nhỏ đính đầy lông vũ lòe loẹt, rồi với ánh mắt tinh quái, cô nhẹ nhàng đặt nó lên đầu Aki.

``A\dots\ Cái này\dots''

Matsuri phụt cười, suýt văng cả nước miếng: ``Trời đất ơi, đúng kiểu quý cô thượng lưu luôn!''

Mel che miệng cười: ``Aki-chan trông sang chảnh quá đi mất!''

Nàng Dị giới đỏ mặt, tỏ vẻ bối rối: ``S-Sao tự nhiên tôi cũng bị kéo vào chuyện này vậy?''

Matsuri: cười gian Bị kéo vào rồi thì cứ tận hưởng đi chị~!

Cả nhóm bắt đầu cười đùa rộn ràng, khiến một vài khách hàng xung quanh quay lại nhìn với ánh mắt ngạc nhiên. Một số người còn mỉm cười thích thú trước màn trêu chọc vui vẻ này, đôi lúc làm họ nhớ về thời học trò, khi mà họ đã từng làm những trò đùa vô tư ấy.

``Được rồi, Botan-chan, em thắng!'' nàng Lễ hội vẫn ôm bụng cười, nhưng dường như cô nàng đã chịu thua trước nước đi cao tay của nàng Sư tử

Mel thở dài sau những trận cười không ngớt đó: ``Không ngờ em lại biết cách `trả đũa' kiểu này đấy\dots''

Aki thì gỡ chiếc mũ ra, cười trừ về phía Botan: ``Lần sau chị sẽ cẩn thận hơn\dots''

Đáp lại, cô nàng Thế hệ thứ Năm chỉ nhún vai, cười nhẹ: ``Em chỉ muốn công bằng thôi mà.''

Dù gì thì cả nhóm cũng đã có một khoảng thời gian vui vẻ. Sau khi ổn định lại, họ quay lại tập trung vào việc chọn đồ trang trí nghiêm túc hơn, nhưng không ai có thể phủ nhận rằng đây là một trong những khoảnh khắc đáng nhớ nhất của buổi mua sắm này.

\ct

Sau khoảng ba tiếng rưỡi mua sắm trong trung tâm thương mại - mà phần lớn thời gian là để bàn bạc, hoặc nếu không bàn bạc nghiêm túc thì cũng là châm chọc nhau khiến không ít khách hàng lẫn nhân viên ba phần ngán ngẩm, bảy phần bất lực - cuối cùng, họ cũng hoàn tất công đoạn mua sắm.

Và đây mới chỉ là bước thứ hai trong kế hoạch chuẩn bị bữa tiệc sinh nhật bất ngờ cho cậu Tổng của họ.

Sora, A-chan và Mio, khi nhìn vào tờ hóa đơn dài như sớ táo quân mà cả ba mươi người vừa gom lại, đã khá ngạc nhiên. Không phải vì tổng tiền - vì đã lường trước được rằng con số sẽ không hề nhỏ vì - mà vì trong danh sách có những món chẳng hề liên quan gì đến một bữa tiệc sinh nhật.

Nàng Đeo kính đưa tờ hóa đơn lên trước mặt Lamy, rồi chỉ vào hai, ba danh mục gần nhau, nghiêm giọng hỏi: ``Yukihana-san? Phiền cô giải thích tại sao lại có rượu whisky với rượu mạnh trong này không?''

Nàng Yêu tinh tuyết cười khì: ``Tại em muốn thử xem chúng sẽ có vị như thế nào\dots''

Haato đứng ra giải thích: ``Em đã bảo em ấy không được mua rượu rồi, nhưng rồi sau đó em lấy lăn ra ăn vạ đấy ạ\dots\ Mà chị nói sao cơ!?''

A-chan chỉ vào danh mục của ba loại rượu có trong tờ hóa đơn. Nàng Song tính thấy vậy, trừng mắt về phía cô đàn em nghiện rượu kia: ``Lamy!! Chị đã bảo là em chỉ lấy một loại thôi mà!''

``Nhưng mỗi một loại thì chán òm à\~{}'' Lamy đáp lại. ``Em muốn biết rượu ở Etsunan khác như thế nào so với rượu sake á\~{}''

Câu trả lời của cô khiến cho Haachama nóng mặt, nhưng chẳng thể làm gì được.

A-chan thở dài, hết thuốc chữa với cô ấy. Cô quay sang Roboco, rồi hỏi với tông giọng tương tự: ``Còn Roboco-san nữa? Trong này có tới bảy, tám loại cà ri khác nhau đây?''

Nàng Người máy cười trừ: ``Em cũng muốn thử xem chúng khác với cà ri em hay ăn thế nào\dots''

``Em không ngờ chị cũng PON nặng đến thế đó\dots'' Haato tặc lưỡi, bất mãn.

``\textbf{Chị không có PON!!!}'' Roboco đáp lại theo thói quen.

Nàng Đeo kính nghe lời giải thích từ hai người, cũng chỉ xoa thái dương bất lực: ``Lần sau muốn thử thì tự bỏ tiền túi ra mua ấy!''

``Vâng\~{}'' cả hai cô nàng đáp lại, không tỏ vẻ hối lỗi gì.

A-chan tiếp tục, quay ngoắt về nàng Hầu gái: ``Minato-san đâu?''

``V-Vâng!''

``Sao cô lại mua cả đống hương liệu với bột mì thế này?'' cô chỉ vào số lượng nguyên liệu làm bánh được chất đầy trên một cái xe.

Aqua đảo mắt một vòng, rồi chầm chậm quay đầu về phía cặp OkaKoro, những người đang huýt sáo đầy giả tạo, tiếp đó liếc sang Shion và Ayame, những kẻ đang làm bộ trò chuyện một cách đáng ngờ.

``Chuyện này\dots\ em cũng không biết nữa\dots'' cô hoang mang, mắt quay như chong chóng.

Cô thật sự chẳng hiểu sao mình lại ôm về cả đống nguyên liệu thế này. Lúc đầu cô chỉ định mua chút bột mì và nguyên vật liệu, nhưng rồi khi gặp lại cặp đôi Chó-Mèo, thấy giỏ hàng họ đầy ắp các loại hương liệu và gia vị. Tiếp đó, khi gặp Shion và Ayame, họ lại thao thao bất tuyệt về sự quan trọng của hương liệu trong bánh ngọt mà chất đống phụ gia.

Nàng Phù thủy cười một cách đắc ý: ``Hương vị phải đa dạng thì bánh mới ngon chứ!''

``Đúng đó! Bánh sinh nhật mà chỉ có một hương vị thì nhàm lắm, yo!'' nàng Quỷ sừng cười gian.

Aqua toát mồ hôi hột, hoảng loạn cực độ: ``Nhưng mà\dots\ nhiều thế này thì ai sẽ làm hết chỗ bánh đó!?''

Nàng Khuyển vỗ vai vào người tiền bối của mình: ``Không sao đâu, cứ tin vào bọn em đi\~{}''

Nàng Miêu cười nhẹ phụ họa: ``Nếu dư thì ta vẫn còn có thể làm bánh khác sau này mà.''

``\dots Nhưng vấn đề là ai sẽ làm!?''

Bốn người chỉ bật cười, chẳng thèm trả lời, càng làm cho Aqua rối bời hơn.

Quay lại hiện tại, Aqua đang cố gắng giải thích nhưng A-chan đã giơ tay ngăn lại. ``Thôi, thôi được rồi, tôi hiểu\dots''

Không biết tại sao, nhưng nàng Hầu gái thở phào nhẹ nhõm vì ít nhất cô không bị phạt.

Sora ở bên cạnh cô cười an ủi: ``Có vẻ như ai cũng có những món đồ ``đặc biệt'' của mình nhỉ\dots''

``Đúng là như vậy đấy,'' A-chan tiếp lời, tiếp tục chỉ vào một món đồ kỳ quặc nữa trong tờ hóa đơn. ``Shirakami-san và Oozora-san?''

Nàng Vịt và nàng Cáo trắng đồng thanh: ``Dạ?''

A-chan cầm tờ hóa đơn lên, mắt lướt qua một danh mục dài ngoằng toàn rau củ quả, rồi liếc sang hai cô nàng trước mặt. ``Đây là của hai em hả?''

Subaru chỉ vào Sora và Choco, đáp: ``Không chỉ bọn em đâu, Sora-senpai với Choco-sensei cũng chọn nữa!''

``Thật sao?''

Nàng Y tá gật đầu: ``Bọn em không biết loại nào hợp với món ăn trong tiệc sinh nhật, nên cứ lấy những loại rau mà mình biết thôi.''

Nàng Thần tượng nhìn lại số rau củ mà họ đã mua chất trên xe, cười với sự ái ngại.

Nàng Đeo kính thở dài, nhìn vào danh sách. Cô thấy rau muống, rau cải, rau xà lách, cà chua, cà rốt, củ cải trắng, bí đỏ, hành tây, thậm chí cả bắp cải tím. Đúng là tươi, sạch, nhưng cũng quá phong phú.

``Tôi không nghĩ bữa tiệc này cần đến nhiều rau củ thế này đâu,'' nàng Đeo kính bất lực.

``Nhưng ít nhất nó sẽ tốt cho sức khỏe mà! Không như ai đó ghét rau củ\dots'' cô nhoẻn mắt về phía Suisei, người không hề để ý tới những gì mà họ đang nói.

A-chan tạm gác vấn đề đó ra một bên, rồi tiếp tục xem hóa đơn: ``Được rồi\dots'' rồi bất chợt, nheo mắt lại đầy ngỡ ngàng: ``Đây là cái gì? Mười\dots\ không, tận hai mươi quả dưa hấu là sao!?''

Cả đám nhìn nhau, rồi Fubuki nhanh chóng giơ tay lên. Cô tươi cười đáp: ``Dưa hấu là món ăn mùa hè mà! Ai cũng thích đúng không?''

Mọi người gật gù đồng tình. Dưa hấu đúng là rất hợp với không khí mùa hè, nhất là trong những bữa tiệc lớn.

Dù vậy, nàng ``Tăng ca'' vẫn không khỏi nghi ngờ: ``Nhưng tận gần hai chục quả thì hơi nhiều đó\dots''

Nàng Cáo trắn chớp mắt vô tội: ``Em chỉ muốn chừa lại vài quả để chơi đập dưa thôi mà!''

``\dots Gì cơ?'' A-chan đơ người.

Subaru thở dài: ``Em đã bảo chị ấy là đừng lấy dư, nhưng mà Fubuki-senpai cứ khăng khăng lấy thêm vài quả để chơi đập dưa á, shuba!''

``Vậy là mấy đứa định biến bữa tiệc sinh nhật thành một lễ hội bãi biển à?'' nàng Đeo kính xoa thái dương, ngớ người với ý tưởng điên rồ này.

``Cũng đâu có tệ đâu nhỉ\~{}''

A-chan thật sự không còn gì để nói với con Cáo trắng này nữa. Nhóm của Fubuki và Subaru đúng là vẫn theo phong cách thường ngày, lúc nào cũng có gì đó khác biệt. Nhưng ít nhất, dưa hấu vẫn có thể ăn được, rau củ cũng vẫn có ích cho bữa tiệc\dots\ Chỉ là cách họ chọn hơi ``hoang dã'' mà thôi.

Chưa kịp nghỉ ngơi, A-chan liếc xuống hóa đơn và nhận ra vài thứ còn kỳ lạ hơn cả rau củ lẫn dưa hấu.

``Tại sao lại có nhiều thịt thế này\dots? Usada-san đâu?''

``Có em đây, peko\dots'' nàng Thỏ rón rén tiến về A-chan.

``Giải thích cho tôi tại sao trong danh sách này lại có một lượng thịt bò khổng lồ thế này đi.''

Nàng Thỏ chỉ về phía nàng Hiệp sĩ: ``Là do Noel mua hết đó, peko!''

``Được rồi, Shirogane-san,'' cô quay sang Noel, ``cô nói đi-''

Nhưng chưa kịp nói hết câu, Flare chợt cắt ngang lời cô: ``Trước khi hỏi chuyện đó, bọn em còn phát hiện ra một điều điên rồ hơn nữa!''

``Hả?''

Rushia tiếp lời, mặt vẫn còn hơi xanh xao: ``Siêu thị này\dots\ còn bán cả bò sống-nanodesu!!''

Toàn bộ ba mươi cô gái rơi vào im lặng trong vài giây. Một cơn rùng mình nhẹ lan ra giữa họ trước thông tin này.

``Cái\dots\ cái gì cơ!?'' Mio há hốc mồm kinh ngạc.

``Chịu khó tìm một chút là thấy ngay! Một con bò to tướng, đúng kiểu bò nguyên con, còn sống nhăn răng, đang được rao bán trong một góc của siêu thị!'' Noel liến thoắng đáp, đôi mắt cô sáng rực.

A-chan ôm trán, bị choáng: ``Từ từ\dots\ Tôi không biết phải phản ứng thế nào với cái thông tin này luôn\dots''

``Tiếc là bọn em không có ý định mua nó,'' Flare theo sau cười trừ.

``Tôi nghĩ hai bà bị thần kinh thật rồi, peko.''

``Vậy thì tốt\dots'' nàng Đeo kính thở dài nhẹ nhõm, khi ít nhất họ không bị mất trí thật sự.

``Nhưng mà Noel-senpai đã mua gần hết thịt bò trong siêu thị rồi, peko!'' nàng Thỏ rít lên đầy bất mãn.

``\dots Gì?''

Pekora không hề ngần ngại mà chỉ vào những danh mục trong tờ hóa đơn, dài bằng một gang tay: ``Đây này! Bà ấy đã mua cả sườn bò, đùi bò, thăn bò nội, thăn bò ngoại\dots\ gần như tất cả những gì có thể lấy được, peko!''

Cả Sora và A-chan hoàn toàn đơ người. Cô nhìn chằm chằm vào Noel, rồi cúi xuống tờ hóa đơn, xác nhận lại một lần nữa. Và đúng như vậy, danh sách dài dằng dặc toàn thịt bò, với số lượng khủng khiếp đến mức gần như vét sạch hai trên ba gian hàng thịt của siêu thị.

``\dots Có nhất thiết phải mua nhiều thế không?''

Noel cười tươi rói: ``Chắc chắn rồi! Một bữa tiệc không thể thiếu thịt bò được!''

``Nhưng mà vét sạch thịt bò trong siêu thị như vậy có hơi quá không?'' nàng Sói đen lên tiếng, e dè hỏi.

``Nhưng vẫn còn một gian hàng nữa mà?'' nàng Hiệp sĩ chớp mắt, tỏ vẻ ngây thơ. Vừa dứt lời, tất cả mọi người đồng loạt nhìn cô với ánh mắt như muốn thốt lên: ``Thật không thể tin được!''

Sora ho nhẹ một tiếng, phá vỡ sự im lặng đầy ngỡ ngàng của mọi người: ``Nhưng mọi người cũng biết là chúng ta phải có một chế độ ăn hợp lý chứ, đúng không? Chúng ta là idol mà.''

Nàng Thỏ gật đầu cái rụp, hoàn toàn đồng ý với cô: ``Đúng đó, peko! Như em đây này, em ăn rất nhiều cà rốt, peko! Chả bù cho Suisei-senpai lười ăn rau thật sự.''

Ngay lập tức, không khí như đông cứng lại. Khuôn mặt của nàng Ca sĩ tối sầm lại, cả người tỏa ra một luồng sát khí rõ rệt.

Cô nàng cất giọng đầy ``vui vẻ,'' nở một nụ cười đầy đáng sợ: ``Có nhất thiết phải chỉ điểm chị như vậy không?''

Tất cả mọi người (trừ Sora và A-chan) đều khiếp đảm trước ánh mắt chết chóc của Cô ca sĩ thiên tài. Ngay cả Noel, người vừa mới phấn khích với thịt bò, cũng phải chùn bước.

``B-Bình tĩnh nào, Suisei\dots'' Mio tiến tới, vội vàng vỗ vai xoa dịu cơn giận của Suisei.

Cô bạn thân của cô - Miko - cũng cười gượng: ``Đúng rồi đó, Sui-chan! Pekora-chan chỉ đùa thôi mà!''

``Thật không\dots?''

``EEeek!!''

Nàng Đeo kính vẫn đang ôm trán, thở dài: ``Được rồi, được rồi, mọi người bình tĩnh lại\dots\ Tôi vẫn chưa xong đâu, còn vài điều nữa đấy.''

A-chan hắng giọng một cái và tiếp tục hỏi: ``Momosuzu-san và Tsunomaki-san? Sao hai đứa mua nhiều vòng thế?''

Cả nàng ``Hải cẩu'' và nàng Cừu đồng thanh hô to: ``Bọn em muốn lập một bộ sưu tập ạ!''

Cả nhóm sững người trước câu trả lời hồn nhiên của hai cô gái vô tư lự này.

``Có\dots\ nhất thiết phải lắm thế không? Ở chỗ ta có nhiều chỗ bán mà?'' nàng ``Tăng ca'' sững người lại với lý do trên.

Watame ôm túi vòng như một báu vật: ``Nhưng ở đây có nhiều loại mà ta không có, A-chan à!''

Đôi mắt Nene thì lấp lánh, rõ ràng tỏ vẻ thích thú với chúng: ``Với lại em mua chúng để khoe với các `chồng' nữa chứ!''

``\dots Hả?'' mắt của A-chan giật giật. Sora chỉ mỉm cười đáp lại: ``Ý của em ấy là các fan hâm mộ của em ấy đó.''

Nàng Hải cẩu ra vẻ tự hào với những vòng tay xinh xắn này: ``Đúng mà! Các chồng của em chắc chắn sẽ thích lắm!''

``Một thần tượng mà không có vòng đẹp thì sao có thể tỏa sáng được, đúng không?'' nàng Cừu phụ họa.

Nghe hai câu trả lời của họ, A-chan chỉ có thể thở dài một hơi.  Cô định tra hỏi tiếp, nhưng ngay lúc đó, nàng Cừu liền cất giọng hát bài ``Watame không có tội'' một cách đầy cảm xúc, trong khi Nene thì giương đôi mắt long lanh, làm nũng đầy đáng yêu.

Tới đây, cô không thể nói thêm một câu nào nữa. Cô hoàn toàn bị hạ gục.

``\dots Thôi được rồi. Coi như tôi chưa thấy gì.''

Toàn bộ nhóm bật cười trước phản ứng bất lực của cô nàng trợ lý.

``Được rồi, tiếp theo\dots\ Rosenthal-san, Natsuiro-san, Yozora-san, Shishiro-san, mấy đứa mua nhiều bờm tóc vậy làm gì?''

Nàng Lễ hội cười khúc khích: ``Để tặng mọi người chứ sao!''

``Nhưng mà ta tới đây để mua đồ phục vụ cho sinh nhật của Hikawa-san mà?''

Nàng Dơi tiếp lời Matsuri: ``Thì đúng là vậy, nhưng ngoài việc mua cho anh ấy, bọn em cũng muốn tặng cho nhau một chút kỷ niệm nữa!''

Nàng Tóc bay vẫy vẫy một chiếc bờm hươu lên, cười một cách thích thú: ``Với cả, nhìn chúng cũng dễ thương lắm á!''

Nàng Sư tử trắng nhún vai: ``Với cả, thêm một món quà cho mọi người thì cũng có sao đâu? Chúng cũng rẻ mà.''

Nhìn những chiếc bờm bắt mắt, nàng Đeo kính cũng đành bó tay, chỉ ngán ngẩm nói: ``\dots Thôi thì, coi như đó là lý do hợp lý.''

Ngay lúc đó, nhóm bốn người họ liền lấy ra một đống bờm tóc với đủ kiểu dáng - từ tai mèo, tai thỏ, cho tới những mẫu kỳ lạ như sừng quỷ hay tai cáo. Họ bắt đầu phân phát cho mọi người, và chẳng mấy chốc, cả nhóm đã biến thành một ``bầy thú'' đầy màu sắc.

Miko đội tai mèo, rên lên thích thú: ``Nya\~{} Nya\~{} Vui quá đi\~{}''

Fubuki cười lớn, đeo chiếc bờm tuần lộc: ``Oi oi, đội cái này vào nhìn ai cũng dễ thương ghê!''

A-chan thở dài, quay về ba cô gái còn lại của Thế hệ thứ Tư: ``Còn Tokomiya-san? Amane-san? Himemori-san? Tôi thấy ba người cũng mua cả đống vòng cổ đấy?''

Towa chống nạnh đáp, cười khì: ``Là quà cho mọi người thôi! Cũng giống như Nene-chan và Watame đó!''

Kanata giơ ra một chiếc vòng có cánh nhỏ, hớn hở: ``Với lại, cái này hợp với chị lắm nè, A-chan!''

``\dots Ừm, nhìn cũng đẹp thật,'' nàng Đeo kính lẩm bẩm, đeo vào cổ tay của mình.

Luna hí hửng phân phát những vòng tay xinh xin tới mọi người: ``Na\~{} Mọi ngừi cũng đeo cái nì đi na\~{}!''

Không khí trở nên vui vẻ hơn hẳn khi mọi người nhận được vòng và bờm tóc, thử đeo chúng và chụp ảnh cùng nhau. Chẳng mấy chốc, một khu của siêu thị trông không khác gì một khu chụp ảnh cosplay (đóng vai) tự sướng, khiến cho những người đi ngang qua không khỏi trầm trồ và thích thú.

``Thôi nào, còn một món nữa, để tôi hỏi nốt đây,'' A-chan quay trở về với sự nghiêm túc vốn có, rồi chỉ vào một món hàng ở gần cuối tờ hóa đơn: ``Được rồi, vậy còn số đá khô này là sao?''

Nàng Vu nữ cười đầy hớn hở: ``Cái đó là để trang trí căn phòng cho thêm phần hoành tráng đó, ne!''

``\textbf{\dots Cái gì cơ!?}'' nàng Sói đen nghe vậy thảng thốt.

``Em không nhớ nói sao? Chị muốn trang trí căn phòng trông thật hoành tráng, ne\~{}''

``\textbf{Tất nhiên là nhớ rồi!} Em tưởng là em đã cản chị rồi mà!?''

``Nhưng mà chị vẫn thấy nó hợp lắm ne\~{}'' Miko cười tít mắt, càng làm cho Mio tức hơn.

A-chan thấy vậy, thở dài chán chường: ``Nói cách khác là cô đã lén mua nó đúng không?''

Nàng Vu nữ vẫn đang nở nụ cười mãn nguyện, giả ngơ không để ý.

Mio đưa hóa đơn lên xem, rồi liếc về thùng giữ nhiệt đang lờ mờ bốc khói ấy, khẽ run rẩy: ``C-Chắc là vẫn trả lại được\dots\ đúng không?''

Nàng Đeo kính liếc về một dòng chữ nhỏ ở dưới cùng của tờ hóa đơn, rồi khẽ lắc đầu. Cái lắc đầu phủ định này khiến cho nàng Sói đen đứng hình.

``Ể? Em tưởng là ta có thể đổi trả hàng chứ?'' Aki lên tiếng, tỏ vẻ hơi ngạc nhiên.

``Đúng là như thế thật, nhưng mà trong này có ghi,'' A-chan đáp lại, rồi liếc về tờ hóa đơn, đọc to dòng chữ đó lên: ``Sản phẩm thuộc nhóm hóa chất đặc biệt, đã mua thì không thể hoàn trả.''

``\dots Tức là bọn em sẽ phải lấy nó về ạ?'' Mio run rẩy, như thể không thể tin vào những gì cô vừa nghe.

Miko nghe vậy, càng phấn khích hơn, thậm chí còn cất giọng như đang trêu chọc Mio: ``Đúng rồi á, ne\~{}''

Nàng Sói đen tới đây cay cú, chỉ có thể siết chặt tay lại mà không thể làm được gì. Cô hít một hơi thật sâu, rồi thở dài bất lực: ``\dots Đành vậy chứ làm sao nữa.''

Cả nhóm không nhịn nổi mà cười khúc khích trước sự cam chịu của nàng Sói. Mio thì chỉ còn biết nhắm mắt lắc đầu, tự nhủ rằng mình sẽ phải giám sát Miko thật chặt khi dùng chỗ đá khô này.

\ct

A-chan vỗ tay, kết thúc buổi tra hỏi: ``Được rồi, vậy là tất cả đều rõ ràng rồi nhé! Giờ thì đến lúc chuyển sang bước tiếp theo!''

Sora quay đầu nhìn mọi người, hơi lưỡng lự một chút nhưng rồi cũng mỉm cười đầy tự tin: ``M-Mọi người à, minh nghĩ là chúng ta coi như đã thành công bước đầu rồi.''

``Vậy chúng ta sẽ làm gì tiếp theo đây, Sora-senpai'' Kanata giơ tay hỏi. Những người khác cũng bắt đầu xĩ xào, ai cũng tò mò về bước tiếp theo trong kế hoạch bí mật này.

Và rồi, nàng Thần tượng kỳ cựu nắm chặt tay, tuyên bố hùng hồn với một nụ cười tỏa sáng: ``\textbf{Chúng ta sẽ đến nhà của Kuon-senpai!}''

Toàn bộ nhóm người chợt im lặng vài giây, trước khi Subaru nhướng mày: ``Khoan đã\dots\ chị vừa nói cái gì cơ?''

``Đến nhà anh ấy bây giờ luôn!''  Matsuri hồ hởi.

``Bình tĩnh nào mọi người! Chúng ta đâu thể xông vào đó ngay được!?'' nàng Sói đen cố gắng phân trần. ``Kuon-senpai sẽ phát hiện ra mất!''

``Tất nhiên là chúng ta chưa chạm trán với anh ấy ngay đâu!'' Sora bật cười. ``Chúng ta đến để chuẩn bị trong bí mật mà!''

Tất cả mọi người thở phào nhẹ nhõm khi được ``Trưởng đoàn'' đính chính lại.

Fubuki cười gian: ``Đúng rồi ha\~{} Đây là tiệc sinh nhật bất ngờ cho anh ấy mừ\~{}''

``Em cứ tưởng là bọn mình sẽ xông vào đột kích chứ\dots'' Nene cười khổ, để rồi nhận lấy một cái cốc đầu nhẹ từ Towa.

``Làm thế chắc anh ấy đau tim luôn á,'' Aqua nhăn mặt, khẽ lắc đầu ngán ngẩm.

A-chan bật cười, vỗ tay thêm lần nữa để tập hợp lại lực lượng: "``óm lại là bây giờ chúng ta sẽ tập hợp lại, chia nhóm hợp lý để chuẩn bị phần tiếp theo của kế hoạch. Không ai được manh động, không ai được tự ý hành động ngoài dự tính, nghe rõ chưa?''

Cả ba mươi người cùng đồng thanh: ``\textbf{Rõ, thưa sếp!}''

Và thế là, sau những giờ phút mua sắm đầy hỗn loạn, mọi người cuối cùng cũng đã hoàn tất việc chuẩn bị nguyên liệu, đồ ăn vặt và quà tặng phục vụ cho buổi tiệc sinh nhật cho Kuon. Kế hoạch đang tiến triển tốt đẹp\dots\ nhưng điều đó không có nghĩa là những bất ngờ phía trước đã kết thúc.

\end{document}